% Mechanically extracted after \mainmatter from the frozen standalone master.

\markboth{Exposé XIX. Reductive Groups---Generalities}{Exposé XIX. Reductive Groups---Generalities}
% BEGIN inputs/XIX/tex/components/00_expose_XIX_title_intro_and_section1_opening_page001.tex
% Authority: Exp19-13oct24.pdf, local p.1 / combined p.1035 / printed p.637.

\chapter*{EXPOSÉ XIX\\[0.4em]
REDUCTIVE GROUPS---GENERALITIES}
\addcontentsline{toc}{chapter}{Exposé XIX. Reductive Groups---Generalities}

\begin{center}
\textit{by M. Demazure}
\end{center}

\begingroup
\renewcommand{\thefootnote}{0}%
\footnotetext{Editor's note: version of 13 October 2024.}%
\endgroup

The remainder of this Seminar is devoted to the study of reductive groups.  Its
principal aim is to generalize Chevalley's classical results (the
\emph{Bible} and \emph{Tôhoku})%
\footnote{Editor's note: see the bibliographical references at the end of this
Exposé.  In particular, the 2005 re-edition of the Chevalley Seminar
1956--1958, cited as \emph{Bible}, was revised by P.\ Cartier; see [Ch05].}
to arbitrary base schemes.  The two central results are the uniqueness theorem
(Exposé XXIII, Theorem 4.1 and \SGARefLink{sga3:xix:paragraph:5-1}{Corollaries 5.1}--5.10) and the existence theorem
(\SGARefLink{sga3:xxv:theorem:1:1}{Exposé XXV, Theorem 1.1}) for pinned reductive group schemes corresponding to
prescribed root data.  The proofs used here are inspired by Chevalley's; the
language of schemes makes it possible to give them greater effectiveness.

The results of the first volume of the \emph{Bible} (Exposés 1--13) will be
used systematically.  By contrast, we shall prove directly over an arbitrary
scheme the results of the second volume (in particular the isomorphism
theorem); knowledge of the proofs over an algebraically closed field is
therefore not absolutely indispensable.

In proving these two fundamental results, we shall use only the most elementary
results from the theory of groups of multiplicative type, found essentially in
Exposés VIII and IX; on the other hand, we shall make essential use of the
results of Exposé XVIII.%
\footnote{Editor's note: more precisely, Proposition \SGARefLink{sga3:xviii:prop:2.3}{XVIII 2.3} (extension of a
``generic homomorphism'' between groups) is used in \SGARefLink{sga3:xxii:prop:4.1.11}{XXII 4.1.11} and then in the
proof of uniqueness in XXIII, as well as in XXIV.  Theorem \SGARefLink{sga3:xviii:thm:3.7}{XVIII 3.7}
(construction of a group from a group germ) is used only in XXV.}
The reader interested specifically in the existence and uniqueness theorems
may, on a first reading, skip Exposés X--XVII.

\SGASection{1}{Review of groups over an algebraically closed field}
\SGARefTarget{sga3:xix:section:1}

\noindent\SGAInlineRefTarget{sga3:xix:paragraph:1-1}\textbf{1.1.}\quad
Throughout this subsection, \(k\) will denote an algebraically closed field.
As announced above, the only results from the \emph{Bible} used below are
contained in its first volume (Exposés 1--13).

% END inputs/XIX/tex/components/00_expose_XIX_title_intro_and_section1_opening_page001.tex

% BEGIN inputs/XIX/tex/components/01_expose_XIX_section1_1_through_1_3_page002.tex
% Authority: Exp19-13oct24.pdf, local p.2 / combined p.1036 / printed pp.638--639.

All the results of the \emph{Bible} (\emph{loc.\ cit.}) concerning semisimple
groups hold more generally for reductive groups (defined below), and their
proofs are identical after the following harmless modifications:
\begin{itemize}
\item Exposé 9, §4, Definition 3: see 1.6.1 below;
\item Exposé 12, §4, Theorem 2(e): omit ``finite'';
\item Exposé 13, §3, Theorem 2: replace ``rank'' by ``semisimple rank'';
\item Exposé 13, §4, Corollary 2 to Theorem 3: replace ``rank'' by
  ``reductive rank''.
\end{itemize}

\noindent\SGAInlineRefTarget{sga3:xix:paragraph:1-2}\textbf{1.2.}\quad
Let \(G\) be a smooth affine connected \(k\)-group.  The \emph{radical} of
\(G\) (\emph{Bible}, §9.4, Proposition 2)%
\footnote{Editor's note: the reference 9-06 (that is, Exposé 9, p.\ 6) has
been replaced by §9.4 (that is, Exposé 9, §4), which applies both to the
\emph{Bible} and to [Ch05].  When the numbering in [Ch05] differs from that
in the \emph{Bible}, as happens in Exposé 6, both references will be given
explicitly.}
is the reduced subgroup associated with the identity component of the
intersection of the Borel subgroups of \(G\); it is also the largest smooth
connected solvable normal subgroup of \(G\).  We denote it by
\(\operatorname{rad}(G)\).

The \emph{unipotent radical} of \(G\) is the unipotent part of the radical of
\(G\); it is also the largest smooth connected unipotent normal subgroup of
\(G\).  We denote it by \(\operatorname{rad}^{u}(G)\).

\noindent\SGAInlineRefTarget{sga3:xix:paragraph:1-3}\textbf{1.3.}\quad
Let \(G\) be a smooth affine connected \(k\)-group, and let \(Q\) be a torus
in \(G\).  Then the centralizer
\(\underline{\operatorname{Centr}}_{G}(Q)\) of \(Q\) in \(G\) is a closed
subgroup of \(G\) (\SGARefLink{sga3:viii:corollary:6:7}{Exposé VIII, 6.7}), smooth (\SGARefLink{sga3:xi:corollary:2:4}{Exposé XI, 2.4}), and connected
(cf.\ \emph{Bible}, §6.6, Theorem 6(a), or [Ch05], §6.7, Theorem 7(a)).

We have the fundamental relation
\begin{equation}
  \operatorname{rad}^{u}\!\bigl(\underline{\operatorname{Centr}}_{G}(Q)\bigr)
  =
  \operatorname{rad}^{u}(G)\cap
  \underline{\operatorname{Centr}}_{G}(Q).
  \tag{1.3.1}\label{sga3:xix:eq:1:3:1}
\end{equation}

First,%
\footnote{Editor's note: the argument that follows has been expanded.}
\[
  U=\operatorname{rad}^{u}(G)\cap
  \underline{\operatorname{Centr}}_{G}(Q)
\]
is a normal unipotent subgroup of
\(\underline{\operatorname{Centr}}_{G}(Q)\); let us show that it is smooth
and connected.  Let \(Q\) act on \(\operatorname{rad}^{u}(G)\) by inner
automorphisms.  Then \(U\) is precisely the scheme of fixed points of this
action.  This scheme is smooth and connected by \SGARefLink{sga3:xix:lemma:1:4}{Lemma 1.4} below.

Consequently, \(U\) is a closed subgroup of
\(\operatorname{rad}^{u}(\underline{\operatorname{Centr}}_{G}(Q))\).
On the other hand, by the \emph{Bible}, Exposé 12, §3, Corollary to
Theorem 1,
\[
  \operatorname{rad}^{u}\!\bigl(
    \underline{\operatorname{Centr}}_{G}(Q)\bigr)(k)
  \subset
  \operatorname{rad}^{u}(G)(k).
\]
Equality \eqref{sga3:xix:eq:1:3:1} follows.

Let us record one special case of the preceding result.  If \(G\) is a smooth
affine connected \(k\)-group and \(T\) is a maximal torus of \(G\), then
\begin{equation}
  \underline{\operatorname{Centr}}_{G}(T)
  =
  T\cdot\bigl(
    \underline{\operatorname{Centr}}_{G}(T)\cap
    \operatorname{rad}^{u}(G)\bigr).
  \tag{1.3.2}\label{sga3:xix:eq:1:3:2}
\end{equation}
Indeed (cf.\ \emph{Bible}, §6.6, Theorem 6(c), or [Ch05], §6.7,
Theorem 7(c)), \(\underline{\operatorname{Centr}}_{G}(T)\) is solvable and
hence is the semidirect product of its maximal torus \(T\) and its unipotent
radical.

By the density theorem (cf.\ \emph{Bible}, §6.5, Theorem 5(a), or [Ch05],
§6.6, Theorem 6(a)), the union of
\(\underline{\operatorname{Centr}}_{G}(T)\), as \(T\) ranges over the
maximal tori of \(G\), contains a dense open subset of \(G\).  It therefore
follows from \eqref{sga3:xix:eq:1:3:2} that:

% END inputs/XIX/tex/components/01_expose_XIX_section1_1_through_1_3_page002.tex

% BEGIN inputs/XIX/tex/components/02_expose_XIX_corollary1_3_3_through_section1_5_page003.tex
% Authority: Exp19-13oct24.pdf, local p.3 / combined p.1037 / printed p.639.

\SGAStatement{Corollary 1.3.3}{%
If \(G\) is a smooth affine connected \(k\)-group, then the union of
\(T\cdot\operatorname{rad}^{u}(G)\), where \(T\) ranges over the maximal tori
of \(G\), contains a dense open subset of \(G\).}
\SGARefTarget{sga3:xix:corollary:1:3:3}

\par\medskip\noindent
\textbf{\textit{\SGARefLink{sga3:xix:notation:1:4:0}{Notation 1.4.0}.}}---%
\footnote{Editor's note: this notation, which appeared in \SGARefLink{sga3:xix:lemma:2-3}{2.3} of the original,
has been inserted here; moreover, the statement and proof of \SGARefLink{sga3:xix:lemma:1-4}{1.4} have been
expanded.}
\textit{We shall systematically use the following notation.  If \(S\) is a
scheme and \(s\) is a point of \(S\), then \(\bar{s}\) denotes the spectrum of
an algebraic closure \(\overline{\kappa(s)}\) of \(\kappa(s)\).}
\SGARefTarget{sga3:xix:notation:1:4:0}

\par\medskip\noindent
\SGAInlineRefTarget{sga3:xix:lemma:1-4}\textbf{\textit{\SGARefLink{sga3:xix:lemma:1-4}{Lemma 1.4}.}}---\footnotemark[5]
\textit{Let \(S\) be a scheme and let \(Q\) be an \(S\)-torus acting on an
\(S\)-group scheme \(H\), separated and smooth over \(S\).
\begin{enumerate}
\item[\SGAInlineRefTarget{sga3:xix:notation:1:4:0:item:i}(i)] The fixed-point functor \(H^{Q}\) is representable by a closed
  subscheme of \(H\), smooth over \(S\).
\item[(ii)] If, moreover, \(H\) is affine over \(S\) and has connected fibers,
  then the same holds for \(H^{Q}\).
\end{enumerate}}
\SGARefTarget{sga3:xix:lemma:1:4}

\SGAProof{Proof}
First, by \SGARefLink{sga3:viii:examples:6:5:item:d:01}{VIII 6.5(d)},%
\footnote{Editor's note: see also the additions made in
VI\textsubscript{B}, \SGARefLink{ex:sga3-VIB-6.2.4-d}{6.2.4(d)} and \SGARefLink{concl:sga3-VIB-6.5.5-a}{6.5.5(a)}.}
since \(H\) is separated over \(S\) and \(Q\) is essentially free over \(S\),
\(H^{Q}\) is representable by a closed subscheme of \(H\).  In particular, if
\(H\) is affine over \(S\), then so is \(H^{Q}\).

Now consider the semidirect product \(G=H\cdot Q\); it is smooth and separated
over \(S\).  The centralizer \(\underline{\operatorname{Centr}}_{G}(Q)\) is
then representable by a closed subscheme of \(G\), of finite presentation over
\(G\), by \SGARefLink{sga3:xi:corollary:6:11:item:a}{Exposé XI 6.11(a)},\footnotemark[6] and it is smooth over \(S\) by
\SGARefLink{sga3:xi:corollary:2:4}{Exposé XI 2.4}.

Since
\[
  \underline{\operatorname{Centr}}_{G}(Q)=H^{Q}\cdot Q,
\]
it follows that \(H^{Q}\) is smooth over \(S\).  Indeed, using the section
\[
  \sigma:H^{Q}\longrightarrow
  \underline{\operatorname{Centr}}_{G}(Q)
\]
deduced from the identity section \(\epsilon_{Q}\) of \(Q\), one sees at once
that \(H^{Q}\) is formally smooth over \(S\); and since \(\epsilon_{Q}\) is of
finite presentation, \(H^{Q}\) is locally of finite presentation over \(S\).
This proves \SGARefLink{sga3:xix:notation:1:4:0:item:i}{(i)}.

Finally, suppose that \(H\) is affine over \(S\) and has connected fibers.
Then every geometric fiber \(G_{\bar{s}}\) of \(G\) is a smooth affine
connected \(\kappa(\bar{s})\)-group.  Hence, by the first assertion of \SGARefLink{sga3:xix:paragraph:1-3}{1.3},
the group
\[
  \underline{\operatorname{Centr}}_{G_{\bar{s}}}(Q_{\bar{s}})
  =
  \bigl(\underline{\operatorname{Centr}}_{G}(Q)\bigr)_{\bar{s}}
  =
  (H^{Q})_{\bar{s}}\cdot Q_{\bar{s}}
\]
is also connected, and this implies that \((H^{Q})_{\bar{s}}\) is connected.
\hfill\(\square\)

\par\medskip\noindent\SGAInlineRefTarget{sga3:xix:paragraph:1-5}\textbf{1.5.}\quad
Recall that the \emph{reductive rank} of a smooth affine \(k\)-group \(G\) is
the common dimension of its maximal tori.  We denote it by
\(\operatorname{rgred}(G/k)\) or \(\operatorname{rgred}(G)\).  The equality
\(\operatorname{rgred}(G/k)=0\) holds if and only if \(G\) is unipotent
(that is, \(G=\operatorname{rad}^{u}(G)\)), by the \emph{Bible}, §6.4,
Corollary 1 to \SGARefLink{sga3:xix:lemma:1:4}{Theorem 4}, or [Ch05], §6.5, Corollary 1 to \SGARefLink{sga3:xix:theorem:2:5}{Theorem 5}.

If \(H\) is a normal subgroup of the smooth affine \(k\)-group \(G\), then the
quotient \(G/H\) is affine and smooth (Exposé VI\textsubscript{B}, \SGARefLink{thm:sga3-VIB-11.17}{11.17} and
VI\textsubscript{A}, \SGARefLink{thm:sga3-VIA-3.3.2:item:iii}{3.3.2(iii)}).  Moreover (\emph{Bible}, §7.3,
Theorem 3(a),(c)),
\[
  \operatorname{rgred}(G)
  =
  \operatorname{rgred}(G/H)+\operatorname{rgred}(H).
\]

% END inputs/XIX/tex/components/02_expose_XIX_corollary1_3_3_through_section1_5_page003.tex

% BEGIN inputs/XIX/tex/components/03_expose_XIX_definition1_6_1_through_section1_7_page004.tex
\par\medskip\noindent
\SGAInlineRefTarget{sga3:xix:definition:1-6-1}\textbf{\textit{\SGARefLink{sga3:xix:definition:1:6:1}{Definition 1.6.1}.}}---\footnote{Editor's note: the former
number 1.6 has been replaced by numbers \SGARefLink{sga3:xix:definition:1-6-1}{1.6.1} through \SGARefLink{sga3:xix:remark:1-6-3}{1.6.3}.  Note that in
this Seminar every reductive \(k\)-group (respectively, every semisimple
\(k\)-group; cf.\ \SGARefLink{sga3:xix:paragraph:1-8}{1.8}) is connected by definition.}
\textit{A \(k\)-group \(G\) is said to be \emph{reductive} if it is smooth,
affine, and connected, and if \(\operatorname{rad}(G)\) is a torus; that is,
if
\[
  \operatorname{rad}^{u}(G)=\{e\}.
\]}
\SGARefTarget{sga3:xix:definition:1:6:1}

\par\medskip\noindent
\textbf{\textit{\SGARefLink{sga3:xix:lemma:1:6:2}{Lemma 1.6.2}.}}---
\textit{Let \(G\) be a reductive \(k\)-group.
\begin{enumerate}
\item[\SGAInlineRefTarget{sga3:xix:definition:1:6:1:item:i}(i)]\SGAInlineRefTarget{sga3:xix:lemma:1:6:2:item:i} If \(Q\) is a torus of \(G\), then
  \(\underline{\operatorname{Centr}}_{G}(Q)\) is reductive.
\item[\SGAInlineRefTarget{sga3:xix:definition:1:6:1:item:ii}(ii)]\SGAInlineRefTarget{sga3:xix:lemma:1:6:2:item:ii} In particular, if \(T\) is a maximal torus of \(G\), then
  \(\underline{\operatorname{Centr}}_{G}(T)=T\).
\item[\SGAInlineRefTarget{sga3:xix:definition:1:6:1:item:iii}(iii)] The center of \(G\) is contained in every maximal torus and is
  therefore diagonalizable.
\item[\SGAInlineRefTarget{sga3:xix:definition:1:6:1:item:iv}(iv)]\SGAInlineRefTarget{sga3:xix:lemma:1:6:2:item:iv} The radical of \(G\) is the unique maximal torus of
  \(\underline{\operatorname{Centr}}(G)\).
\end{enumerate}}
\SGARefTarget{sga3:xix:lemma:1:6:2}

\SGAProof{Proof}
Indeed, \SGARefLink{sga3:xix:definition:1:6:1:item:i}{(i)} follows from (\SGARefLink{sga3:xix:eq:1:3:1}{1.3.1}), \SGARefLink{sga3:xix:definition:1:6:1:item:ii}{(ii)} from (\SGARefLink{sga3:xix:eq:1:3:2}{1.3.2}), and \SGARefLink{sga3:xix:definition:1:6:1:item:iii}{(iii)} follows from
\SGARefLink{sga3:xix:definition:1:6:1:item:ii}{(ii)}.  Finally, the maximal torus of
\(\underline{\operatorname{Centr}}(G)\), that is, the neutral component
\(\underline{\operatorname{Centr}}(G)^{0}\), is a smooth connected solvable
subgroup normal in \(G\), and hence is contained in
\(\operatorname{rad}(G)\).  Conversely, \(\operatorname{rad}(G)\) is a torus
normal in \(G\), and hence central (\emph{Bible}, §\SGARefLink{sga3:xix:paragraph:4-3}{4.3}, Corollary to
Proposition 2); this proves \SGARefLink{sga3:xix:definition:1:6:1:item:iv}{(iv)}.
\hfill\(\square\)

\par\medskip\noindent
\SGAInlineRefTarget{sga3:xix:remark:1-6-3}\textbf{\textit{\SGARefLink{sga3:xix:remark:1:6:3}{Remark 1.6.3}.}}---
\textit{If \(G\) is reductive and
\(\dim(G)\ne\operatorname{rgred}(G)\), then
\[
  \dim(G)-\operatorname{rgred}(G)\geq 2.
\]
Indeed, this difference is always even (cf.\ \SGARefLink{sga3:xix:paragraph:1-10}{1.10} below).}
\SGARefTarget{sga3:xix:remark:1:6:3}

\par\medskip\noindent\SGAInlineRefTarget{sga3:xix:paragraph:1-7}\textbf{1.7.}\quad
Let \(G\) be a smooth affine connected \(k\)-group, and let \(H\) be a
smooth connected normal subgroup.  Then
\begin{equation}
  \operatorname{rad}(H)
  =
  \bigl(\operatorname{rad}(G)\cap H\bigr)_{\mathrm{red}}^{0},
  \qquad
  \operatorname{rad}^{u}(H)
  =
  \bigl(\operatorname{rad}^{u}(G)\cap H\bigr)_{\mathrm{red}}^{0},
  \tag{1.7.1}
\end{equation}
as is seen immediately.  In particular, if \(G\) is reductive, then so is
\(H\).
\SGARefTarget{sga3:xix:equation:1:7:1}

Let \(f:G\longrightarrow G'\) be a surjective morphism of smooth affine
connected \(k\)-groups.  Then
\begin{equation}
  f\bigl(\operatorname{rad}^{u}(G)\bigr)
  =
  \operatorname{rad}^{u}(G').
  \tag{1.7.2}
\end{equation}
In particular, if \(G\) is reductive, then so is \(G'\).
\SGARefTarget{sga3:xix:equation:1:7:2}

Let us prove (\SGARefLink{sga3:xix:equation:1:7:2}{1.7.2}).  First,
\(f\bigl(\operatorname{rad}^{u}(G)\bigr)\) is contained in
\(\operatorname{rad}^{u}(G')\).  On introducing
\[
  H=
  \bigl(f^{-1}(\operatorname{rad}^{u}(G'))\bigr)_{\mathrm{red}}^{0},
\]
which contains \(\operatorname{rad}^{u}(G)\), we have
\(\operatorname{rad}^{u}(H)=\operatorname{rad}^{u}(G)\); thus we are reduced
to the case \(G=H\), that is, to the case in which \(G'\) is unipotent.
As \(T\) runs over the maximal tori of \(G\), the union of the groups
\(T\operatorname{rad}^{u}(G)\) is dense in \(G\).  Hence the union of the
groups \(f(T)\,f\bigl(\operatorname{rad}^{u}(G)\bigr)\) is dense in \(G'\).
But \(f(T)\) consists of semisimple elements, so \(f(T)=\{e\}\), since
\(G'\) is unipotent.  It follows that
\(f\bigl(\operatorname{rad}^{u}(G)\bigr)\) is dense in \(G'\).  Therefore,
by \emph{Bible}, §\SGARefLink{sga3:xix:paragraph:5:4}{5.4}, \SGARefLink{sga3:xix:lemma:1:4}{Lemma 4}, or [Ch05], §6.1, \SGARefLink{sga3:xix:lemma:3:1}{Lemma 1},\footnote{Editor's
note: the following argument has been supplied in detail.}
\(f\bigl(\operatorname{rad}^{u}(G)\bigr)\) is an open subgroup of \(G'\).
Since \(G'\) is connected, it follows that
\[
  f\bigl(\operatorname{rad}^{u}(G)\bigr)=G'.
\]

\noindent\textit{N.B.}\quad Under the same hypotheses one can prove that
\[
  f\bigl(\operatorname{rad}(G)\bigr)=\operatorname{rad}(G').
\]

% END inputs/XIX/tex/components/03_expose_XIX_definition1_6_1_through_section1_7_page004.tex

% BEGIN inputs/XIX/tex/components/04_expose_XIX_sections1_8_through_1_10_pages005_006.tex
\par\medskip\noindent\SGAInlineRefTarget{sga3:xix:paragraph:1-8}\textbf{1.8.}\quad
A \(k\)-group \(G\) is said to be \emph{semisimple} if it is smooth,
affine, and connected, and if
\[
  \operatorname{rad}(G)=\{e\},
\]
that is, if \(G\) is reductive and
\(\underline{\operatorname{Centr}}(G)\) is finite.  If \(G\) is an
arbitrary smooth affine connected \(k\)-group, then
\(G/\operatorname{rad}(G)\) is semisimple (\emph{Bible}, §9.4,
Proposition 2), and \(G/\operatorname{rad}^{u}(G)\) is reductive.  The
\emph{semisimple rank of \(G\)}, denoted by
\(\operatorname{rgss}(G/k)\), or simply by \(\operatorname{rgss}(G)\), is
the reductive rank of \(G/\operatorname{rad}(G)\).
\SGARefTarget{sga3:xix:section:1:8}

Thus, if \(G\) is reductive,
\[
  \operatorname{rgred}(G)
  =
  \operatorname{rgss}(G)+\dim\bigl(\operatorname{rad}(G)\bigr).
\]

If \(G\) is a smooth affine connected \(k\)-group and \(Q\) is a central
subtorus of \(G\), then \(G/Q\) is semisimple if and only if \(G\) is
reductive and \(Q=\operatorname{rad}(G)\).  Indeed, if \(G/Q\) is
semisimple, then \(Q\supset\operatorname{rad}(G)\); hence
\(\operatorname{rad}(G)\) is a torus, so \(G\) is reductive, and
\(\operatorname{rad}(G)\) is the maximal torus of
\(\underline{\operatorname{Centr}}(G)\), whence
\(\operatorname{rad}(G)=Q\).  If \(G\) is reductive and \(Q\) is a central
subgroup, then \(G/Q\) is reductive and
\[
  \operatorname{rgss}(G)=\operatorname{rgss}(G/Q).
\]

\par\medskip\noindent\SGAInlineRefTarget{sga3:xix:paragraph:1-9}\textbf{1.9.}\quad
Let \(K\) be an algebraically closed extension of \(k\), and let \(G\) be
a smooth affine connected \(k\)-group.  Then \(G\) is reductive
(respectively, semisimple) if and only if \(G_K\) is, and
\[
  \operatorname{rgred}(G/k)=\operatorname{rgred}(G_K/K),
  \qquad
  \operatorname{rgss}(G/k)=\operatorname{rgss}(G_K/K).
\]
\SGARefTarget{sga3:xix:section:1:9}

\par\medskip\noindent\SGAInlineRefTarget{sga3:xix:paragraph:1-10}\textbf{1.10.}\quad
Let \(G\) be a smooth connected \(k\)-group and let \(T\) be a torus of
\(G\).\footnote{Editor's note: the following sentence has been added.}
Write \(\mathfrak g\) for the \(k\)-Lie algebra of \(G\), that is,
\(\mathfrak g=\operatorname{Lie}(G)=\operatorname{Lie}(G/k)(k)\), and
similarly write \(\mathfrak t=\operatorname{Lie}(T)\).  Under the action
of \(T\), through \(T\longrightarrow G\) and the adjoint representation
of \(G\), the Lie algebra \(\mathfrak g\) decomposes as
\begin{equation}
  \mathfrak g
  =
  \mathfrak g^{0}
  \oplus
  \bigoplus_{\alpha\in R}\mathfrak g^{\alpha},
  \tag{1.10.1}\label{sga3:xix:paragraph:1-10:equation:1-10-1}
\end{equation}
where the elements \(\alpha\in R\) are nontrivial characters of \(T\)
and the spaces \(\mathfrak g^{\alpha}\) are nonzero.  The characters
\(\alpha\in R\) are called the \emph{roots of \(G\) relative to \(T\)}.
By \SGARefLink{sga3:ii:theorem:5-2-3:item:i}{Exposé II, 5.2.3(i)},
\begin{equation}
  \mathfrak g^{0}
  =
  \operatorname{Lie}\bigl(\underline{\operatorname{Centr}}_{G}(T)\bigr).
  \tag{1.10.2}
\end{equation}
\SGARefTarget{sga3:xix:equation:1:10:1}
\SGARefTarget{sga3:xix:equation:1:10:2}

In particular,\footnote{Editor's note: the following argument has been
given in detail.} since
\(\underline{\operatorname{Centr}}_{G}(T)\) is connected (by \SGARefLink{sga3:xix:paragraph:1-3}{1.3} when
\(G\) is affine, and by \SGARefLink{sga3:xii:6.6:item:b}{Exposé XII, 6.6(b)}, in general), \(T\) is its own
centralizer if and only if
\(\mathfrak g^{0}=\operatorname{Lie}(T)\).  This condition implies that
\(T\) is maximal and that
\(\underline{\operatorname{Centr}}(G)\subset T\); hence, by \SGARefLink{sga3:xii:8.8:item:d:02}{Exposé XII,
8.8(d)},
\begin{equation}
  \underline{\operatorname{Centr}}(G)
  =
  \operatorname{Ker}\bigl(T\longrightarrow
    \operatorname{GL}(\mathfrak g)\bigr)
  =
  \bigcap_{\alpha\in R}\operatorname{Ker}(\alpha).
  \tag{1.10.3}
\end{equation}
\SGARefTarget{sga3:xix:equation:1:10:3}
Moreover, \(\underline{\operatorname{Centr}}(G)\) is then affine, and
therefore the morphism
\(G\longrightarrow G/\underline{\operatorname{Centr}}(G)\) is affine.
Since the latter group is affine (\SGARefLink{sga3:xii:6.1}{Exposé XII, 6.1}),\footnote{Editor's
note: see also Addition VI\textsubscript{B}, \SGARefLink{cor:sga3-VIB-6.2.6}{6.2.6}.} \(G\) is affine.

When \(G\) is reductive and \(T\) is maximal, the roots in the preceding
sense coincide with the roots in the sense of \emph{Bible}, §12.2,
Definition 1.  The latter are indeed roots in the sense of this number
(\emph{Bible}, §13.2, Theorem 1(c)), and their number is
\(\dim(G)-\dim(T)\) (\emph{Bible}, §13.4, Corollary 2 to Theorem 3).
Moreover, if \(\alpha\) is a root, then so is \(-\alpha\)
(\emph{Bible}, §12.2, Corollary to Proposition 1).  As usual, we write
the group law of
\(\operatorname{Hom}_{k\text{-}\mathrm{gr}}
  (T,\mathbb G_{m,k})\)
indifferently additively or multiplicatively.  It follows that, for
reductive \(G\),
\begin{equation}
  \dim(G)-\operatorname{rgred}(G)=\operatorname{Card}(R)
  \tag{1.10.4}
\end{equation}
is always even.
\SGARefTarget{sga3:xix:equation:1:10:4}

The semisimple rank of the reductive group \(G\) is the rank of the
subset \(R\) of the \(\mathbb Q\)-vector space
\[
  \operatorname{Hom}_{k\text{-}\mathrm{gr}}
    (T,\mathbb G_{m,k})\otimes_{\mathbb Z}\mathbb Q
\]
(\emph{Bible}, §13.3, Theorem 2).
\SGARefTarget{sga3:xix:section:1:10}

% END inputs/XIX/tex/components/04_expose_XIX_sections1_8_through_1_10_pages005_006.tex

% BEGIN inputs/XIX/tex/components/05_expose_XIX_lemma1_11_opening_page006.tex
\par\medskip\noindent
\SGAInlineRefTarget{sga3:xix:lemma:1-11}\textbf{\textit{\SGARefLink{sga3:xix:lemma:1:11}{Lemma 1.11}.}}---
\textit{Let \(k\) be an algebraically closed field, \(G\) a smooth affine
connected \(k\)-group, and \(T\) a torus of \(G\).  Put
\[
  W(T)
  =
  \underline{\operatorname{Norm}}_{G}(T)/
  \underline{\operatorname{Centr}}_{G}(T),
\]
the Weyl group of \(G\) relative to \(T\).  The following conditions are
equivalent:
\begin{enumerate}
\item[(i)] \(G\) is reductive, \(T\) is maximal, and
  \(\operatorname{rgss}(G)=1\).
\item[(ii)] \(G\) is reductive, \(T\) is maximal, and \(G\ne T\); there
  exists a subtorus \(Q\) of \(T\), of codimension \(1\) in \(T\), that
  is central in \(G\).
\item[(iii)] \(G\) is not solvable and
  \(\dim(G)-\dim(T)\leq 2\).
\item[(iv)] \(W(T)\ne\{e\}\) and
  \(\dim(G)-\dim(T)\leq 2\).
\item[(v)] \(W(T)=(\mathbb Z/2\mathbb Z)_k\) and
  \(\dim(G)-\dim(T)=2\).
\end{enumerate}
Moreover, under these conditions there are exactly two roots of \(G\)
relative to \(T\), and they are opposite.  Under the conditions of
\textup{(ii)}, \(Q=\operatorname{rad}(G)\).}
\SGARefTarget{sga3:xix:lemma:1:11}

\SGAProof{Proof}
The implication (v)\(\Rightarrow\)(iv) is clear, and
(iv)\(\Rightarrow\)\SGARefLink{sga3:xix:lemma:1:11:item:iii}{(iii)} follows from \emph{Bible}, §6.1, Corollary 3 to
Theorem 1, or [Ch05], §6.2, Corollary 3 to Theorem 2.

Let us prove \SGARefLink{sga3:xix:lemma:1:11:item:iii}{(iii)}\(\Rightarrow\)\SGARefLink{sga3:xix:lemma:1:11:item:ii}{(ii)}.  Let \(U\) be the unipotent radical
of \(G\).  The group \(G/U\) is reductive and is not a torus (otherwise
\(G\) would be solvable); hence, by (\SGARefLink{sga3:xix:equation:1:10:4}{1.10.4}),
\[
  \dim(G/U)-\operatorname{rgred}(G/U)
  =
  \operatorname{Card}(R)
  \geq 2.
\]
On the other hand,
\[
  \operatorname{rgred}(G)
  =
  \operatorname{rgred}(G/U)+\operatorname{rgred}(U)
  =
  \operatorname{rgred}(G/U),
\]
and therefore
\[
\begin{aligned}
  \dim(G)-\dim(U)-\operatorname{rgred}(G)
    &=\dim(G/U)-\operatorname{rgred}(G/U)\\
    &\geq 2
     \geq \dim(G)-\dim(T)
     \geq \dim(G)-\operatorname{rgred}(G).
\end{aligned}
\]
It follows that \(\dim(U)=0\), so \(G\) is reductive, and that
\(\operatorname{rgred}(G)=\dim(T)\); thus \(T\) is maximal and
\(\dim(G)-\dim(T)=2\).  By \emph{Bible}, §10.4, Proposition 8, there
exists a singular torus \(Q\) of codimension \(1\) in \(T\).  Then
\(\underline{\operatorname{Centr}}_{G}(Q)\) is reductive (1.6.2(i)) and
not solvable (by the definition of a singular torus), and hence
\[
  \dim\bigl(\underline{\operatorname{Centr}}_{G}(Q)\bigr)
  -\operatorname{rgred}(G)\geq 2.
\]
This proves that
\(G=\underline{\operatorname{Centr}}_{G}(Q)\), and that \(Q\) is central
in \(G\).

Let us prove \SGARefLink{sga3:xix:lemma:1:11:item:ii}{(ii)}\(\Rightarrow\)\SGARefLink{sga3:xix:lemma:1:11:item:i}{(i)}.  We know that \(G/Q\) is reductive
by \SGARefLink{sga3:xix:paragraph:1-7}{1.7} and that \(\operatorname{rgred}(G/Q)=1\); hence
\(\operatorname{rgss}(G/Q)\) is either \(0\) or \(1\).  The first case is
impossible, since it would imply \(\operatorname{rgss}(G)=0\), and
therefore \(G=T\).  Thus
\[
  \operatorname{rgred}(G/Q)=\operatorname{rgss}(G/Q)=1,
\]
which proves that \(G/Q\) is semisimple.  Consequently \(Q\) is the
radical of \(G\) and \(\operatorname{rgss}(G)=1\).

% END inputs/XIX/tex/components/05_expose_XIX_lemma1_11_opening_page006.tex

% BEGIN inputs/XIX/tex/components/06_expose_XIX_lemma1_11_completion_and_proposition1_12_pages007_008.tex
Let us finally prove (i)\(\Rightarrow\)(v).  If \(Q\) is the radical of
\(G\), then
\(\dim(T)-\dim(Q)=1\), and \(Q\) is central in \(G\).  Hence
\(G=\underline{\operatorname{Centr}}_{G}(Q)\), which proves that \(Q\)
is a singular torus.  By \emph{Bible}, §11.3, Theorem 2,
\[
  W_G(T)=(\mathbb Z/2\mathbb Z)_k,
\]
and, by \emph{Bible}, §12.1, Lemma 1,
\(\dim(G)-\dim(T)=2\).  Thus \(G\) has at most two roots relative to
\(T\), while it has at least two, opposite to one another, by \SGARefLink{sga3:xix:paragraph:1-10}{1.10}.

\par\medskip\noindent
\SGAInlineRefTarget{sga3:xix:proposition:1-12}\textbf{\textit{\SGARefLink{sga3:xix:proposition:1:12}{Proposition 1.12}.}}---
\textit{Let \(k\) be an algebraically closed field, \(G\) a smooth
connected \(k\)-group, \(T\) a torus of \(G\), and \(R\) the set of
roots of \(G\) relative to \(T\).  Write
\[
  \mathfrak g
  =
  \mathfrak g^{0}
  \oplus
  \bigoplus_{\alpha\in R}\mathfrak g^{\alpha},
  \qquad
  \mathfrak g^{\alpha}\ne0,
\]
for the decomposition of \(\mathfrak g\) under
\(\operatorname{Ad}(T)\).  For every \(\alpha\in R\), put
\[
  T_{\alpha}=\operatorname{Ker}(\alpha)^{0},
  \qquad
  Z_{\alpha}=\underline{\operatorname{Centr}}_{G}(T_{\alpha}).
\]
\footnote{Editor's note: thus \(T_{\alpha}\) is the identity component
of \(\operatorname{Ker}(\alpha)\).}
The following conditions are equivalent:
\begin{enumerate}
\item[\SGAInlineRefTarget{sga3:xix:lemma:1:11:item:i}(i)] \(G\) is affine and reductive, and \(T\) is maximal.
\item[\SGAInlineRefTarget{sga3:xix:lemma:1:11:item:ii}(ii)] \(\mathfrak g^{0}=\mathfrak t\), and every \(Z_{\alpha}\),
  for \(\alpha\in R\), is reductive.
\item[\SGAInlineRefTarget{sga3:xix:lemma:1:11:item:iii}(iii)] \(\mathfrak g^{0}=\mathfrak t\); every
  \(\mathfrak g^{\alpha}\), for \(\alpha\in R\), has dimension \(1\);
  if \(\alpha,\beta\in R\) and \(\beta=q\alpha\) for some
  \(q\in\mathbb Q\), then \(q=\pm1\); and, for every \(\alpha\in R\),
  there exists \(w_{\alpha}\in G(k)\) which normalizes \(T\),
  centralizes \(T_{\alpha}\), but does not centralize \(T\).
\end{enumerate}
Moreover, under these conditions every \(Z_{\alpha}\) has semisimple
rank \(1\), and
\[
  \operatorname{Lie}(Z_{\alpha})
  =
  \mathfrak t\oplus
  \mathfrak g^{\alpha}\oplus
  \mathfrak g^{-\alpha}.
\]}
\SGARefTarget{sga3:xix:proposition:1:12}

\SGAProof{Proof}
Suppose first that \(G\) is affine and reductive and that \(T\) is
maximal.  Every \(Z_{\alpha}\) is then affine and reductive by \SGARefLink{sga3:xix:lemma:1:6:2:item:i}{1.6.2(i)},
with maximal torus \(T\); furthermore, \(T\) is its own centralizer by
\SGARefLink{sga3:xix:lemma:1:6:2:item:ii}{1.6.2(ii)}.  Thus \(\mathfrak g^{0}=\mathfrak t\), proving
\SGARefLink{sga3:xix:proposition:1:12:item:i}{(i)}\(\Rightarrow\)\SGARefLink{sga3:xix:proposition:1:12:item:ii}{(ii)}.

Conversely, the equality \(\mathfrak g^{0}=\mathfrak t\) implies that
\(T\) is maximal and that \(G\) is affine, by \SGARefLink{sga3:xix:paragraph:1-10}{1.10}.  Hence every
\(Z_{\alpha}\) is affine, smooth, and connected, by \SGARefLink{sga3:xix:paragraph:1-3}{1.3}.  In all cases,
\SGARefLink{sga3:ii:lemma:5-2-2}{Exposé II, 5.2.2} gives
\[
  \operatorname{Lie}(Z_{\alpha})
  =
  \mathfrak g^{T_{\alpha}}
  =
  \mathfrak g^{0}
  \oplus
  \bigoplus_{\beta\in R\cap\mathbb Q\alpha}
    \mathfrak g^{\beta}.
\]
Consequently,
\[
  \operatorname{Lie}(Z_{\alpha})
  \supset
  \mathfrak t\oplus
  \mathfrak g^{\alpha}\oplus
  \mathfrak g^{-\alpha},
\]
so \(Z_{\alpha}\ne T\).  Since \(T_{\alpha}\) is a central subtorus of
\(Z_{\alpha}\), of codimension \(1\) in \(T\), \SGARefLink{sga3:xix:lemma:1-11}{Lemma 1.11} applied to
\(Z_{\alpha}\) proves the equivalence of \SGARefLink{sga3:xix:proposition:1:12:item:ii}{(ii)} and (iii), together with
the additional assertions.

It remains to prove \SGARefLink{sga3:xix:proposition:1:12:item:ii}{(ii)}\(\Rightarrow\)\SGARefLink{sga3:xix:proposition:1:12:item:i}{(i)}.  We already know that \SGARefLink{sga3:xix:proposition:1:12:item:ii}{(ii)}
implies that \(T\) is maximal and that \(G\) is affine.  Let \(U\) be
the unipotent radical of \(G\).  It is invariant in \(G\), and its Lie
algebra is invariant under \(\operatorname{Ad}(T)\); hence
\[
  \operatorname{Lie}(U)
  =
  \bigl(\operatorname{Lie}(U)\cap\mathfrak g^{0}\bigr)
  \oplus
  \bigoplus_{\alpha\in R}
    \bigl(\operatorname{Lie}(U)\cap\mathfrak g^{\alpha}\bigr).
\]
By (\SGARefLink{sga3:xix:eq:1:3:1}{1.3.1}),
\[
\begin{aligned}
  U\cap T
    &=
    U\cap\underline{\operatorname{Centr}}_{G}(T)
     =
    \operatorname{rad}^{u}
      \bigl(\underline{\operatorname{Centr}}_{G}(T)\bigr)
     =
    \operatorname{rad}^{u}(T)
     =
    \{e\},\\
  U\cap Z_{\alpha}
    &=
    U\cap\underline{\operatorname{Centr}}_{G}(T_{\alpha})
     =
    \operatorname{rad}^{u}
      \bigl(\underline{\operatorname{Centr}}_{G}(T_{\alpha})\bigr)
     =
    \operatorname{rad}^{u}(Z_{\alpha})
     =
    \{e\}.
\end{aligned}
\]
It follows\footnote{Editor's note: if \(U\) and \(L\) are two subgroup
schemes of \(G\), then
\(\operatorname{Lie}(U)\cap\operatorname{Lie}(L)
=\operatorname{Lie}(U\cap L)\).} that
\[
\begin{aligned}
  \operatorname{Lie}(U)\cap\mathfrak g^{0}
    &=
    \operatorname{Lie}(U)\cap\mathfrak t
     =
    \operatorname{Lie}(U\cap T)
     =
    0,\\
  \operatorname{Lie}(U)\cap\mathfrak g^{\alpha}
    &\subset
    \operatorname{Lie}(U)\cap\operatorname{Lie}(Z_{\alpha})
     =
    \operatorname{Lie}(U\cap Z_{\alpha})
     =
    0.
\end{aligned}
\]
Thus \(\operatorname{Lie}(U)=0\), and therefore \(U=\{e\}\); in other
words, \(G\) is reductive.

\par\medskip\noindent
\SGAInlineRefTarget{sga3:xix:corollary:1-13}\textbf{\textit{\SGARefLink{sga3:xix:corollary:1:13}{Corollary 1.13}.}}---
\textit{Let \(k\), \(G\), \(T\), \(R\), and
\[
  \mathfrak g
  =
  \mathfrak g^{0}
  \oplus
  \bigoplus_{\alpha\in R}\mathfrak g^{\alpha}
\]
be as in \SGARefLink{sga3:xix:proposition:1-12}{Proposition 1.12}.  For every \(\alpha\in R\), let
\(T_{\alpha}\) be the maximal torus of \(\operatorname{Ker}(\alpha)\)
and put
\(Z_{\alpha}=\underline{\operatorname{Centr}}_{G}(T_{\alpha})\).
The following conditions are equivalent:
\begin{enumerate}
\item[\SGAInlineRefTarget{sga3:xix:proposition:1:12:item:i}(i)] \(G\) is affine and semisimple, and \(T\) is maximal.
\item[\SGAInlineRefTarget{sga3:xix:proposition:1:12:item:ii}(ii)] \(\mathfrak g^{0}=\mathfrak t\), every \(Z_{\alpha}\) is
  reductive, and
  \(\bigcap_{\alpha\in R}\operatorname{Ker}(\alpha)\) is finite.
\end{enumerate}
\footnote{Editor's note: this follows from the equalities
\(\operatorname{rad}(G)
=\underline{\operatorname{Centr}}(G)^{0}\) and
\(\underline{\operatorname{Centr}}(G)
=\bigcap_{\alpha\in R}\operatorname{Ker}(\alpha)\), established in
\SGARefLink{sga3:xix:lemma:1:6:2:item:iv}{1.6.2(iv)} and (\SGARefLink{sga3:xix:equation:1:10:3}{1.10.3}).}}
\SGARefTarget{sga3:xix:corollary:1:13}

% END inputs/XIX/tex/components/06_expose_XIX_lemma1_11_completion_and_proposition1_12_pages007_008.tex

% BEGIN inputs/XIX/tex/components/07_expose_XIX_section2_scholium2_1_and_2_2_pages008_009.tex
\SGASection{2}{Reductive group schemes: definitions and first properties}
\SGARefTarget{sga3:xix:section:2}

\par\medskip\noindent
\textbf{\textit{\SGARefLink{sga3:xix:scholium:2:1}{Scholium 2.1}.}}---
\textit{For an \(S\)-group scheme \(G\), the following properties are
equivalent:
\begin{enumerate}
\item[\SGAInlineRefTarget{sga3:xix:section:2:item:i}(i)] \(G\) is affine and smooth over \(S\), with connected
  fibers.
\item[\SGAInlineRefTarget{sga3:xix:section:2:item:ii}(ii)] \(G\) is affine and flat over \(S\), of finite presentation
  over \(S\), with geometrically integral fibers.
\end{enumerate}
These properties are stable under base change and are local for the
fpqc topology.}
\SGARefTarget{sga3:xix:scholium:2:1}

\footnote{Editor's note: the following details have been added.}
Indeed, suppose that \SGARefLink{sga3:xix:section:2:item:i}{(i)} holds.  Since \(G\) is affine and smooth over
\(S\), it is of finite presentation over \(S\); and since its fibers
are smooth and connected, they are geometrically integral, by
Exposé VI\textsubscript{A}, \SGARefLink{prop:sga3-VIA-2.4}{2.4}.

Conversely, if \SGARefLink{sga3:xix:section:2:item:ii}{(ii)} holds, the fibers of \(G\) are connected and
geometrically reduced, and hence smooth
(Exposé VI\textsubscript{A}, 1.3.1).  It follows that \(G\) is smooth
over \(S\), by EGA IV\textsubscript{4}, 17.5.2.

These properties are plainly stable under base change: see EGA II,
1.5.1 for ``affine'', EGA IV\textsubscript{1}, 1.6.2(iii) for ``of
finite presentation'', EGA IV\textsubscript{2}, 2.1.4 for ``flat'',
and EGA IV\textsubscript{2}, 4.6.5(i) for ``with geometrically integral
fibers''.

They are also clearly local for the Zariski topology.  It is therefore
enough to check that, if \(S'\longrightarrow S\) is faithfully flat and
quasi-compact and \(G_{S'}\longrightarrow S'\) has the indicated
properties, then so does \(G\longrightarrow S\).  This follows from
EGA IV\textsubscript{2}, 2.5.1 for ``flat'', 2.7.1(vi) and (xiii) for
``of finite presentation'' and ``affine'', and 4.6.5(i) for ``with
geometrically integral fibers''; for ``smooth'', see also
EGA IV\textsubscript{4}, 17.7.3(ii).

\par\medskip\noindent\SGAInlineRefTarget{sga3:xix:paragraph:2-2}\textbf{2.2.}\quad
Let \(G\) be an \(S\)-group scheme satisfying the preceding conditions,
and let \(Q\) be a torus of \(G\) in the sense of \SGARefLink{sga3:ix:definition:1:3}{Exposé IX,
Definition 1.3}.\footnote{Editor's note: the following argument has been
given in detail.}
By \SGARefLink{sga3:xi:corollary:6:11:item:a}{Exposé XI, 6.11(a)} and \SGARefLink{sga3:xix:remark:2-4}{2.4},
\(\underline{\operatorname{Centr}}_{G}(Q)\) is represented by a closed
subgroup scheme of \(G\); it is therefore affine over \(S\), and it is
of finite presentation and smooth over \(S\).  Moreover, every
geometric fiber \(G_{\bar s}\) is a smooth affine connected
\(\kappa(\bar s)\)-group.  The first assertion of \SGARefLink{sga3:xix:paragraph:1-3}{1.3} therefore shows
that the same is true of
\[
  \underline{\operatorname{Centr}}_{G_{\bar s}}(Q_{\bar s})
  =
  \bigl(\underline{\operatorname{Centr}}_{G}(Q)\bigr)_{\bar s}.
\]
\SGARefTarget{sga3:xix:section:2:2}

% END inputs/XIX/tex/components/07_expose_XIX_section2_scholium2_1_and_2_2_pages008_009.tex

% BEGIN inputs/XIX/tex/components/08_expose_XIX_lemma2_3_and_remark2_4_page009.tex
\par\medskip\noindent
\SGAInlineRefTarget{sga3:xix:lemma:2-3}\textbf{\textit{\SGARefLink{sga3:xix:lemma:2:3}{Lemma 2.3}.}}---
\textit{Let \(S\) be a scheme, let \(G\) be an \(S\)-group scheme that
is smooth and affine over \(S\), with connected fibers, and let \(T\)
be a torus of \(G\).  The set of points \(s\in S\) for which
\(G_{\bar s}\) is reductive, of semisimple rank \(1\), with maximal
torus \(T_{\bar s}\), is an open subset \(U\) of \(S\).}
\SGARefTarget{sga3:xix:lemma:2:3}

\SGAProof{Proof}
Since \(G\) and \(T\) are flat over \(S\), the function
\[
  s
  \longmapsto
  \dim(G_{\bar s}/\bar s)-\dim(T_{\bar s}/\bar s)
\]
is locally constant.  Let \(U_1\) be the open subset of \(S\) on which
this function has the value \(2\).

\footnote{Editor's note: the following details have been added.  Note
also that §6 is independent of the rest of this Exposé.}
By \SGARefLink{sga3:xix:proposition:6-3}{6.3} the Weyl group
\[
  W_G(T)
  =
  \underline{\operatorname{Norm}}_{G}(T)/
  \underline{\operatorname{Centr}}_{G}(T)
\]
is represented by an \(S\)-group scheme that is étale and separated
over \(S\), and the function
\[
  s\longmapsto
  \text{the number of points of \(W_G(T)_{\bar s}\)}
\]
is lower semicontinuous.  Let \(U_2\) be the locus on which this number
is greater than \(1\); it is open.  By Lemma 1.11, the locus on which
\(G_{\bar s}\) is reductive, of semisimple rank \(1\), with maximal
torus \(T_{\bar s}\), is
\[
  U=U_1\cap U_2.
\]
Furthermore, \(W_G(T)_{\bar s}\) has exactly two points for every
\(s\in U\).  Consequently, by SGA 1, I, 10.9 and
EGA IV\textsubscript{3}, 15.5.1 and IV\textsubscript{4}, 18.12.4,
\(W_G(T)_U\) is finite and étale over \(U\).

\par\medskip\noindent
\SGAInlineRefTarget{sga3:xix:remark:2-4}\textbf{\textit{\SGARefLink{sga3:xix:remark:2:4}{Remark 2.4}.}}---
\textit{The group \(W_G(T)_U\) is isomorphic to
\((\mathbb Z/2\mathbb Z)_U\).}
\SGARefTarget{sga3:xix:remark:2:4}

Indeed, the preceding argument shows that it is finite and étale over
\(U\).  Since the automorphism functor of
\((\mathbb Z/2\mathbb Z)_U\) is the trivial group, it is enough to
verify the assertion locally.  It is immediate whenever the algebra
\(\mathscr A\) defining \(W_G(T)_U\) is a free
\(\mathcal O_U\)-module.\footnote{Editor's note: the hypothesis implies
that \(\mathscr A\) is locally free of rank \(2\).  Since the
augmentation ideal \(\mathscr I\) is a direct summand of \(\mathscr A\),
after replacing \(U\) by a sufficiently small affine open
\(\operatorname{Spec}(R)\), one may suppose that
\(I=\Gamma(U,\mathscr I)\) is a free \(R\)-module of rank \(1\).
If \(x\) is a generator of \(I\), then \(x^2=ax\) for some \(a\in R\);
the assumption that \(A=\Gamma(U,\mathscr A)\) is étale implies that
\(a\) is invertible, which gives the asserted local identification.}

% END inputs/XIX/tex/components/08_expose_XIX_lemma2_3_and_remark2_4_page009.tex

% BEGIN inputs/XIX/tex/components/09_expose_XIX_notations2_5_0_and_theorem2_5_page010.tex
\par\medskip\noindent
\textbf{\textit{Notations and reminders 2.5.0.}}---%
\footnote{Editor's note: this paragraph of notations and reminders has
been added.}
Let \(S\) be a scheme, let \(G\) be an \(S\)-group scheme, and let
\(\epsilon\colon S\to G\) be the unit section of \(G\).  We saw in
II, §\SGARefLink{sga3:xix:paragraph:4-11}{4.11} that the functor
\(\underline{\operatorname{Lie}}(G/S)\) is represented by the vector
bundle
\(\mathbb V(\omega^1_{G/S})\), where
\(\omega^1_{G/S}=\epsilon^{*}(\Omega^1_{G/S})\), and we write
\[
  \mathfrak g
  =
  \operatorname{Lie}(G/S)
  =
  \bigl(\omega^1_{G/S}\bigr)^{\vee}
\]
for the sheaf of sections of this vector bundle.  Suppose in addition
that \(G\) is smooth over \(S\).  Then \(\omega^1_{G/S}\), and hence
\(\operatorname{Lie}(G/S)\), are finite locally free
\(\mathcal O_S\)-modules, and (cf. I, \SGARefLink{sga3:i:corollary:4-6-5-1}{4.6.5.1})
\[
  \underline{\operatorname{Lie}}(G/S)
  =
  \mathbb W\bigl(\operatorname{Lie}(G/S)\bigr);
\]
that is, for every \(S\)-scheme \(S'\),
\[
  \underline{\operatorname{Lie}}(G/S)(S')
  =
  \Gamma\!\left(
    S',
    \operatorname{Lie}(G/S)\otimes_{\mathcal O_S}\mathcal O_{S'}
  \right).
\]
By II, \SGARefLink{sga3:ii:corollary:4-1-1-1}{4.1.1.1}, the adjoint action of \(G\) gives
\(\underline{\operatorname{Lie}}(G/S)
=\mathbb W(\operatorname{Lie}(G/S))\) the structure of a
\(G\)-\(\mathcal O_S\)-module; thus
\(\operatorname{Lie}(G/S)\) is a \(G\)-\(\mathcal O_S\)-module
(cf. I, \SGARefLink{sga3:i:definition:4-7-1}{4.7.1}).  If \(G\) is moreover affine over \(S\), this says,
by I, \SGARefLink{sga3:i:proposition:4-7-2}{4.7.2}, that \(\operatorname{Lie}(G/S)\) is an
\(A(G)\)-comodule.

If \(T\) is a torus over \(S\) (IX, \SGARefLink{sga3:ix:definition:1:3}{Definition 1.3}), we say that \(T\)
is \emph{split} (``trivial'' in the original) if it is isomorphic to
\(\mathbb G_{m,S}^{\,r}\) for some integer \(r\geq 0\), and that \(T\)
is \emph{trivialized} if such an isomorphism has been fixed, or more
generally if an isomorphism
\[
  T\simeq D_S(M)
\]
has been fixed, where \(M\) is a free abelian group of rank \(r\).

Finally, recall (XII, \SGARefLink{sga3:xii:1.3}{Definition 1.3}) that a torus \(T\) of \(G\) is
called \emph{maximal} if, for every \(s\in S\), \(T_{\bar s}\) is a
maximal torus of \(G_{\bar s}\).
\SGARefTarget{sga3:xix:notations:2:5:0}

\par\medskip\noindent
\SGAInlineRefTarget{sga3:xix:theorem:2-5}\textbf{\textit{\SGARefLink{sga3:xix:theorem:2:5}{Theorem 2.5}.}}---
\textit{Let \(S\) be a scheme, let \(G\) be an \(S\)-group scheme that
is affine and of finite presentation over \(S\), with connected
fibers, and let \(s_0\) be a point of \(S\).  Suppose that \(G\) is
flat over \(S\) at \(\epsilon(s_0)\), and that the geometric fiber
\(G_{\bar s_0}\) is (reduced and) reductive
\textup{(resp. semisimple)}.  Then there exist an open subset
\(U\subset S\), containing \(s_0\), and an étale surjective morphism
\(S'\to U\) such that:}
\SGARefTarget{sga3:xix:theorem:2:5}

\begin{enumerate}
\item[\textit{\SGAInlineRefTarget{sga3:xix:theorem:2:5:item:i}(i)}]
\textit{\(G_U\) is smooth, with reductive
\textup{(resp. semisimple)} geometric fibers, and its reductive rank
and semisimple rank are constant.}

\item[\textit{\SGAInlineRefTarget{sga3:xix:theorem:2:5:item:ii}(ii)}]
\textit{\(G_{S'}\) has a split\footnote{Editor's note: the word
``split'' has been added.} maximal torus \(T\), and the Weyl group
\textup{(cf. \SGARefLink{sga3:xix:proposition:6-3}{6.3})}}
\[
  W_{G_{S'}}(T)
  =
  \underline{\operatorname{Norm}}_{G_{S'}}(T)/
  \underline{\operatorname{Centr}}_{G_{S'}}(T)
  =
  \underline{\operatorname{Norm}}_{G_{S'}}(T)/T
\]
\textit{is finite over \(S'\).}
\end{enumerate}


% END inputs/XIX/tex/components/09_expose_XIX_notations2_5_0_and_theorem2_5_page010.tex

% BEGIN inputs/XIX/tex/components/10_expose_XIX_theorem2_5_proof_pages010_012.tex
\SGAProof{Proof}\footnote{Editor's note: the proof has been expanded.}
Write \(\pi\colon G\to S\) for the structure morphism and
\(\epsilon\colon S\to G\) for the unit section.  Since \(G\) is flat
over \(S\) at \(\epsilon(s_0)\), and since \(G_{\bar s_0}\) is reduced
(hence smooth over \(\bar s_0\), by VI\textsubscript{A}, \SGARefLink{prop:sga3-VIA-1.3.1}{1.3.1}),
\(G\) is smooth over \(S\) at \(\epsilon(s_0)\)
(EGA IV\textsubscript{4}, 17.5.1).  Thus there is an open
neighborhood \(V\) of \(\epsilon(s_0)\) on which \(\pi\) is smooth.
Then \(S'=\epsilon^{-1}(V)\) is open in \(S\), and \(G_{S'}\) is
smooth over \(S'\) at the points of \(\epsilon(S')\).  Since \(G\)
has connected fibers, VI\textsubscript{B}, \SGARefLink{thm:sga3-VIB-3.10}{3.10} shows that \(G_{S'}\)
is smooth over \(S'\).  Replacing \(S\) by \(S'\), we may therefore
suppose that \(G\) is smooth over \(S\).

By \SGARefLink{sga3:xi:theorem:4:1}{Exposé XI, Theorem 4.1}, the functor of subgroups of multiplicative
type of \(G\) is represented by an \(S\)-scheme \(\mathscr M\), smooth
and separated over \(S\).  Let \(r_0\) be the rank of
\(G_{\bar s_0}\), and consider the open subscheme
\(\mathscr M_{r_0}\subset\mathscr M\) representing the functor of
subtori of rank \(r_0\) of \(G\) (IX, \SGARefLink{sga3:ix:definition:1:4}{1.4}).  Smoothness implies that
\(\mathscr M_{r_0}\) has a rational point over a finite separable
extension of \(\kappa(s_0)\) (EGA IV\textsubscript{4}, 17.15.10(iii)).
Consequently there is an étale morphism \(S'\to S\), together with a
point \(s'_0\) above \(s_0\), such that \(G_{s'_0}\) has a torus of
rank \(r_0\).  Replacing \(S\) by \(S'\), we may therefore suppose
that \(G_{s_0}\) has a torus of rank \(r_0\).

By ``Hensel's lemma'' (XI, \SGARefLink{sga3:xi:proposition:1:10}{1.10}), the smoothness of
\(\mathscr M_{r_0}\) allows this torus to be lifted to an \(S'\)-torus
\(T\) of \(G\), where \(S'\to S\) is étale and is equipped with a
point \(s'_0\) above \(s_0\).  By \SGARefLink{sga3:x:corollary:4-5}{Exposé X, 4.5}
(see also \SGARefLink{sga3:xix:proposition:6-1}{6.1}\footnote{Editor's note: §\SGARefLink{sga3:xix:proposition:6-1}{6.1} is independent of the rest
of this Exposé.}), there is an étale morphism \(f\colon S''\to S'\)
that splits \(T\) and for which \(f^{-1}(s'_0)\ne\varnothing\).
Because an étale morphism is open, and the assertions in \textup{\SGARefLink{sga3:xix:theorem:2:5:item:i}{(i)}}
are local for the étale topology, we may thus suppose that \(G\) has
a split torus \(T\),\footnote{Editor's note: throughout this volume,
``trivial torus'' has been replaced by ``split torus''.} maximal at
\(s_0\).  Write \(T=D_S(M)\), and let
\[
  \mathfrak g=\prod_{m\in M}\mathfrak g^{m}
\]
be the decomposition of
\(\mathfrak g=\operatorname{Lie}(G/S)\) under
\(\operatorname{Ad}(T)\) (\SGARefLink{sga3:i:proposition:4-7-3}{Exposé I, 4.7.3}).

Put \(\mathfrak t=\operatorname{Lie}(T)\), and, for every \(m\in M\),
write
\[
  \mathfrak g^{m}(s_0)
  =
  \mathfrak g^{m}\otimes_{\mathcal O_S}\kappa(s_0).
\]
\footnote{Editor's note: the preceding sentence has been added.}
Let \(R\) be the set of \(m\in M\), \(m\ne0\), such that
\(\mathfrak g^{m}(s_0)\ne0\).\footnote{Editor's note: since
\(\mathfrak g\) is a finite locally free \(\mathcal O_S\)-module,
\(R\) is finite.}  Since \(G_{\bar s_0}\) is reductive,
\(\mathfrak g^{0}(s_0)=\mathfrak t(s_0)\), and hence
\[
  \mathfrak g(s_0)
  =
  \mathfrak t(s_0)
  \oplus
  \prod_{\alpha\in R}\mathfrak g^{\alpha}(s_0).
\]
Since the modules in question are locally free, after shrinking \(S\)
we may suppose that the \(\mathfrak g^{\alpha}\) are free and that
\[
  \mathfrak g
  =
  \mathfrak t
  \oplus
  \prod_{\alpha\in R}\mathfrak g^{\alpha}.
\]

Recall from \SGARefLink{sga3:xii:1.12}{Exposé XII, 1.12} that a group of multiplicative type has a
unique maximal torus (this also follows almost immediately by descent,
the diagonalizable case being evident).  Let \(T_\alpha\) be the
maximal torus of \(\ker(\alpha)\), and put
\[
  Z_\alpha=\underline{\operatorname{Centr}}_G(T_\alpha).
\]
\footnote{Editor's note: the following argument has been expanded.}
By \SGARefLink{sga3:xix:paragraph:2-2}{2.2}, \(Z_\alpha\) is affine and smooth over \(S\), with connected
fibers; and by \SGARefLink{sga3:xix:proposition:1-12}{1.12}, its fiber
\((Z_\alpha)_{\bar s_0}\) is reductive, of semisimple rank \(1\), with
maximal torus \(T_{\bar s_0}\).  By \SGARefLink{sga3:xix:lemma:2-3}{2.3} there is therefore an open
subset \(U_\alpha\subset S\), containing \(s_0\), on which
\(Z_\alpha\) has reductive fibers.

Set
\[
  U=\bigcap_{\alpha\in R}U_\alpha.
\]
By \SGARefLink{sga3:xix:proposition:1-12}{1.12}, for every \(s\in U\), \(G_{\bar s}\) is reductive, with
maximal torus \(T_{\bar s}\), and the set of roots of \(G_{\bar s}\)
relative to \(T_{\bar s}\) is identified with \(R\).  Consequently
\[
  \operatorname{rgred}(G_{\bar s})
  =
  \dim(T)
  =
  \operatorname{rg}(M),
  \qquad
  \operatorname{rgss}(G_{\bar s})
  =
  \operatorname{rg}(R)
  \quad\text{(cf. 1.10).}
\]
This proves \textup{\SGARefLink{sga3:xix:theorem:2:5:item:i}{(i)}} and the first assertion of \textup{\SGARefLink{sga3:xix:theorem:2:5:item:ii}{(ii)}}.
It remains to prove that \(W_{G_U}(T_U)\) is finite over \(U\), that
is, that it has the same number of points in every geometric fiber
(SGA 1, I, 10.9; EGA IV\textsubscript{3}, 15.5.1; and
EGA IV\textsubscript{4}, 8.12.4).

For this it is enough to observe that, for \(s\in U\), the geometric
fiber of this group is determined by the configuration \(R\subset M\):
it is the constant group associated with the abstract Weyl group of
this root system.  In particular it is independent of \(s\); see the
\emph{Bible}, §11.3, Theorem 2 (and also Exposé XXII, no. 3).
\hfill\(\square\)


% END inputs/XIX/tex/components/10_expose_XIX_theorem2_5_proof_pages010_012.tex

% BEGIN inputs/XIX/tex/components/11_expose_XIX_corollary2_6_through_remark2_9_page012.tex
\par\medskip\noindent
\textbf{\textit{\SGARefLink{sga3:xix:corollary:2:6}{Corollary 2.6}.}}---
\textit{Let \(G\) be an \(S\)-group scheme that is affine and smooth
over \(S\), with connected fibers.  The set of points \(s\in S\) for
which \(G_{\bar s}\) is reductive \textup{(resp. semisimple)} is an
open subset \(U\) of \(S\), and the functions}
\[
  s\longmapsto\operatorname{rgred}(G_{\bar s}/\bar s),
  \qquad
  s\longmapsto\operatorname{rgss}(G_{\bar s}/\bar s)
\]
\textit{are locally constant on \(U\).}
\SGARefTarget{sga3:xix:corollary:2:6}

\par\medskip\noindent
\textbf{\textit{\SGARefLink{sga3:xix:definition:2:7}{Definition 2.7}.}}---
\textit{An \(S\)-group \textup{(that is, an \(S\)-group scheme)}
\(G\) is called \emph{reductive} \textup{(resp. \emph{semisimple})}
if it is affine and smooth over \(S\), with connected reductive
\textup{(resp. semisimple)} geometric fibers.}
\SGARefTarget{sga3:xix:definition:2:7}

For an \(S\)-group scheme \(G\), the property of being reductive
(resp. semisimple) is stable under base change and local for the fpqc
topology.

\par\medskip\noindent
\textbf{2.8.}---
Let \(G\) be a reductive \(S\)-group.  For every torus
(resp. maximal torus) \(Q\) of \(G\),
\[
  Z(Q)=\underline{\operatorname{Centr}}_G(Q)
\]
is reductive (resp. is equal to \(Q\)).  This follows at once from \SGARefLink{sga3:xix:paragraph:2-2}{2.2}
and 1.6.  Applying \SGARefLink{sga3:xix:theorem:2-5}{2.5} to
\(\underline{\operatorname{Centr}}_G(Q)\), we deduce that \(Q\) is
contained, locally for the étale topology, in a maximal torus.
\SGARefTarget{sga3:xix:paragraph:2:8}

\par\medskip\noindent
\textbf{\textit{Remark 2.9.}}---
\textit{Using, as usual, the technique of EGA IV\textsubscript{3},
§8, the finite-presentation hypotheses, and Theorem 2.5, one sees
that if \(G\) is a reductive group over \(S\), there is an open cover
\(\{U_i\}\) of \(S\) such that each \(G_{U_i}\) is obtained by base
change from a reductive group over a noetherian ring \textup{(in fact,
over a finitely generated \(\mathbb Z\)-algebra)}.  Likewise, if \(G\)
has a split maximal torus \(T\), one may further suppose that
\(T_{U_i}\) comes from a split maximal torus of the preceding
group, \ldots}
\SGARefTarget{sga3:xix:remark:2:9}

\SGASection{3}{Roots and root systems of reductive group schemes}
\SGARefTarget{sga3:xix:section:3}


% END inputs/XIX/tex/components/11_expose_XIX_corollary2_6_through_remark2_9_page012.tex

% BEGIN inputs/XIX/tex/components/12_expose_XIX_section3_1_through_3_4_page013.tex
\par\medskip\noindent
\textbf{\SGAInlineRefTarget{sga3:xix:paragraph:3-1}3.1.}---
Let \(S\) be a scheme, and let \(T\) be an \(S\)-torus acting linearly
on a finite locally free \(\mathcal O_S\)-module \(F\) (cf. I, §\SGARefLink{sga3:xix:lemma:4-7}{4.7}).
For every character \(\alpha\) of \(T\), that is,
\[
  \alpha\in
  \operatorname{Hom}_{S\text{-gr.}}(T,\mathbb G_{m,S}),
\]
define a subfunctor of \(\mathbb W(F)\) by
\[
  \mathbb W(F)^\alpha(S')
  =
  \left\{
    x\in\mathbb W(F)(S')
    \ \middle|\
    t\cdot x=\alpha(t)x
    \text{ for every }t\in T(S''),\ S''\to S'
  \right\}.
\]
\SGARefTarget{sga3:xix:paragraph:3:1}

\par\medskip\noindent
\textbf{\textit{Lemma.}}---%
\footnote{Editor's note: this statement has been made into a lemma in
order to bring it out.}
\textit{Then
\(\mathbb W(F)^\alpha=\mathbb W(F^\alpha)\), where \(F^\alpha\) is a
submodule of \(F\), locally a direct summand of \(F\), and hence also
locally free.}
\SGARefTarget{sga3:xix:lemma:3:1}

Indeed, the assertion is local for the fpqc topology, and we may
suppose that \(T=D_S(M)\), where \(M\) is a finite-type (free) abelian
group.  Then \(\alpha\) is identified with a locally constant function
from \(S\) to \(M\) (\SGARefLink{sga3:viii:corollary:1:3}{Exposé VIII, 1.3}), and, after restricting \(S\) if
necessary, we may suppose that this function is constant.  We are then
reduced to \SGARefLink{sga3:i:proposition:4-7-3}{Exposé I, 4.7.3}.

\par\medskip\noindent
\SGAInlineRefTarget{sga3:xix:definition:3-2}\textbf{\textit{\SGARefLink{sga3:xix:definition:3:2}{Definition 3.2}.}}---
\textit{Let \(S\) be a scheme, let \(G\) be an \(S\)-group scheme
smooth over \(S\) and with connected fibers,\footnote{Editor's note:
the original also assumed that \(G\) was of finite presentation over
\(S\), an assumption that does not appear to be used below.  In any
case, since \(G\) is smooth over \(S\) and has connected fibers, it is
quasi-compact and separated over \(S\) (VI\textsubscript{B}, \SGARefLink{cor:sga3-VIB-5.5}{5.5}), and
hence of finite presentation over \(S\).} and let \(T\) be a torus of
\(G\).  Put \(\mathfrak g=\operatorname{Lie}(G/S)\), and let \(T\)
act on \(\mathfrak g\) through the adjoint representation of \(G\).}

\textit{A character \(\alpha\) of \(T\) is called a \emph{root} of
\(G\) with respect to \(T\) if the following equivalent conditions
are satisfied:}
\SGARefTarget{sga3:xix:definition:3:2}

\begin{enumerate}
\item[\textit{(i)}]
\textit{For every \(s\in S\), \(\alpha_s\) is a root of \(G_s\) with
respect to \(T_s\) (\SGARefLink{sga3:xix:paragraph:1-10}{1.10}).}

\item[\textit{(ii)}]
\textit{\(\alpha\) is nontrivial on every fiber, and
\(\mathfrak g^\alpha(s)\ne0\) for every \(s\in S\).}
\end{enumerate}

\par\medskip\noindent
\textbf{\textit{\SGARefLink{sga3:xix:lemma:3:3}{Lemma 3.3}.}}---
\textit{Let \(S\) be a scheme, let \(T\) be an \(S\)-torus, and let
\(\alpha\) be a character of \(T\).  The following conditions are
equivalent:}
\SGARefTarget{sga3:xix:lemma:3:3}

\begin{enumerate}
\item[\textit{\SGAInlineRefTarget{sga3:xix:lemma:3:3:item:i}(i)}]
\textit{\(\alpha\) is nontrivial on every fiber; that is, for every
\(s\in S\), \(\alpha_s\) is distinct from the unit character of
\(T_s\).}

\item[\textit{\SGAInlineRefTarget{sga3:xix:lemma:3:3:item:ii}(ii)}]
\textit{For every \(S'\to S\) with \(S'\ne\varnothing\),
\(\alpha_{S'}\) is distinct from the unit character of \(T_{S'}\).}

\item[\textit{\SGAInlineRefTarget{sga3:xix:lemma:3:3:item:iii}(iii)}]
\textit{The morphism \(\alpha\) is faithfully flat.}
\end{enumerate}

\footnote{Editor's note: the original has been simplified; it invoked
\SGARefLink{sga3:ix:corollary:5:2}{Exposé IX, 5.2}.  Note, however, that loc.\ cit., \SGARefLink{sga3:xix:lemma:5-3}{5.3} shows that, if
\(S\) is connected, the conditions of the lemma hold as soon as
\textup{\SGARefLink{sga3:xix:lemma:3:3:item:i}{(i)}} holds at one point \(s\) of \(S\).}
It is clear that \SGARefLink{sga3:xix:lemma:3:3:item:ii}{(ii)} implies \SGARefLink{sga3:xix:lemma:3:3:item:i}{(i)}, and it is easy to see that \SGARefLink{sga3:xix:lemma:3:3:item:iii}{(iii)}
implies \SGARefLink{sga3:xix:lemma:3:3:item:i}{(i)}.  Condition \SGARefLink{sga3:xix:lemma:3:3:item:i}{(i)} implies \SGARefLink{sga3:xix:lemma:3:3:item:ii}{(ii)}: if \(s'\in S'\) lies over
\(s\), and \(\alpha_{s'}\) is the unit character, then so is
\(\alpha_s\).  Finally, to prove that \SGARefLink{sga3:xix:lemma:3:3:item:i}{(i)} implies \SGARefLink{sga3:xix:lemma:3:3:item:iii}{(iii)}, we reduce to
the case in which \(T\) is diagonalizable and conclude by \SGARefLink{sga3:viii:corollary:3:2:item:a:01}{Exposé VIII,
3.2(a)}.

\par\medskip\noindent
\textbf{3.4.}---%
\footnote{Editor's note: what follows has been modified so as to recall
the hypotheses of \SGARefLink{sga3:xix:definition:3-2}{3.2} and to add that \(T\) is a maximal torus.}
Let \(G\) be a reductive \(S\)-group scheme, let \(T\) be a maximal
torus of \(G\), and let \(\alpha\) be a root of \(G\) with respect to
\(T\).  Then, by \SGARefLink{sga3:xix:notations:2:5:0}{2.5.0} and \SGARefLink{sga3:xix:proposition:1-12}{1.12}, \(\mathfrak g^\alpha\) is a locally
free \(\mathcal O_S\)-module of rank one.  Moreover, by \SGARefLink{sga3:xix:paragraph:1-10}{1.10},
\(-\alpha\) is also a root of \(G\) with respect to \(T\).  In
particular, if \(G\) has semisimple rank one, \SGARefLink{sga3:xix:lemma:1-11}{1.11} and \SGARefLink{sga3:xix:proposition:1-12}{1.12} give the
following statement.
\SGARefTarget{sga3:xix:paragraph:3:4}

% END inputs/XIX/tex/components/12_expose_XIX_section3_1_through_3_4_page013.tex

% BEGIN inputs/XIX/tex/components/13_expose_XIX_lemma3_5_through_root_functor_page014.tex
\par\medskip\noindent
\textbf{\textit{\SGARefLink{sga3:xix:lemma:3:5}{Lemma 3.5}.}}---
\textit{Let \(S\) be a scheme, let \(G\) be a reductive \(S\)-group
of semisimple rank one, and let \(T\) be a maximal torus of \(G\).  If
\(\alpha\) is a root of \(G\) with respect to \(T\), then so is
\(-\alpha\), and}
\[
  \mathfrak g
  =
  \mathfrak t
  \oplus \mathfrak g^\alpha
  \oplus \mathfrak g^{-\alpha},
\]
\textit{where \(\mathfrak g^\alpha\) and
\(\mathfrak g^{-\alpha}\) are locally free of rank one.}
\SGARefTarget{sga3:xix:lemma:3:5}

\par\medskip\noindent
\SGAInlineRefTarget{sga3:xix:definition:3-6}\textbf{\textit{\SGARefLink{sga3:xix:definition:3:6}{Definition 3.6}.}}---
\textit{Let \(S\) be a scheme, let \(G\) be a reductive \(S\)-group,
let \(T\) be a maximal torus of \(G\), and let \(R\) be a set of roots
of \(G\) with respect to \(T\).  We say that \(R\) is a \emph{root
system} of \(G\) with respect to \(T\) if the following equivalent
conditions are satisfied:}
\SGARefTarget{sga3:xix:definition:3:6}

\begin{enumerate}
\item[\textit{(i)}]
\textit{For every \(s\in S\), the map
\(\alpha\mapsto\alpha_s\) is a bijection from \(R\) onto the set of
roots of \(G_s\) with respect to \(T_s\).}

\item[\textit{\SGAInlineRefTarget{sga3:xix:definition:3:6:item:ii}(ii)}]
\textit{The elements of \(R\) are distinct on every fiber
\textup{(that is, if \(\alpha,\alpha'\in R\) and
\(\alpha\ne\alpha'\), then
\(\alpha-\alpha'=\alpha{\alpha'}^{-1}\) is nontrivial on every
fiber)}, and, for every \(s\in S\),}
\[
  \dim(G_s/\kappa(s))-\dim(T_s/\kappa(s))=\operatorname{Card}(R).
\]

\item[\textit{\SGAInlineRefTarget{sga3:xix:definition:3:6:item:iii}(iii)}]
\textit{One has}
\[
  \mathfrak g
  =
  \mathfrak t
  \oplus
  \bigoplus_{\alpha\in R}\mathfrak g^\alpha.
\]
\end{enumerate}

The equivalence of these conditions is immediate.

\par\medskip\noindent
\SGAInlineRefTarget{sga3:xix:lemma:3-7}\textbf{\textit{\SGARefLink{sga3:xix:lemma:3:7}{Lemma 3.7}.}}---
\textit{Let \(S\) be a scheme, let \(G\) be a reductive \(S\)-group,
let \(T\) be a maximal torus of \(G\), and let \(R\) be a root system
of \(G\) with respect to \(T\).  Every root of \(G\) with respect to
\(T\) is locally on \(S\) equal to an element of \(R\).}
\SGARefTarget{sga3:xix:lemma:3:7}

\SGAProof{Proof}
This follows from condition \SGARefLink{sga3:xix:definition:3:6:item:iii}{(iii)}.

Put
\[
  \mathcal M
  =
  \underline{\operatorname{Hom}}_{S\text{-gr.}}
  (T,\mathbb G_{m,S});
\]
this is a twisted constant \(S\)-group scheme (\SGARefLink{sga3:x:theorem:5-6}{Exposé X, 5.6}).  If
\(G\) has a root system \(R\) with respect to \(T\), then the inclusion
\(R\hookrightarrow\mathcal M(S)\) defines a morphism
\[
  R_S\longrightarrow\mathcal M,
\]
where \(R_S\) is the constant \(S\)-scheme defined by \(R\).  By
condition \SGARefLink{sga3:xix:definition:3:6:item:ii}{(ii)}, this morphism is an open and closed immersion whose
image is precisely
\[
  \bigcup_{\alpha\in R}\alpha(S),
\]
where each \(\alpha\in R\) is regarded as a section
\(S\to\mathcal M\).

Let \(\mathcal R\) be the functor of roots of \(G\) with respect to
\(T\).  By definition, \(\mathcal R(S')\) is the set of roots of
\(G_{S'}\) with respect to \(T_{S'}\), for every \(S'\to S\).  If
\(S'=\varnothing\), set \(\mathcal R(\varnothing)=\{e\}\); if
\(S'\ne\varnothing\), the inclusion
\(R\hookrightarrow\mathcal M(S')\) identifies \(R\) with a root
system of \(G_{S'}\) with respect to \(T_{S'}\).  Thus, by \SGARefLink{sga3:xix:lemma:3-7}{3.7},
\[
  \mathcal R(S')
  =
  \operatorname{Hom}_{\mathrm{loc.\ const.}}(S',R),
\]
which shows that \(\mathcal R\) is represented by \(R_S\).

Now drop the assumption that \(G\) necessarily has a root system with
respect to \(T\).  The functor \(\mathcal R\) is in any case an fpqc
subsheaf of \(\mathcal M\).  Locally for this topology, \(G\) does
have a root system with respect to \(T\) (for example, make \(T\)
split and use the proof of \SGARefLink{sga3:xix:theorem:2:5}{Theorem 2.5}).  By \SGARefLink{sga3:iv:proposition:4:6:8}{Exposé IV, 4.6.8} and the
descent theory of open (respectively closed)
subschemes,\footnote{Editor's note: see SGA 1, VIII, 4.4 or
EGA IV\textsubscript{2}, 2.7.1.} we obtain the following proposition.

% END inputs/XIX/tex/components/13_expose_XIX_lemma3_5_through_root_functor_page014.tex

% BEGIN inputs/XIX/tex/components/14_expose_XIX_proposition3_8_through_3_10_page015.tex
\par\medskip\noindent
\SGAInlineRefTarget{sga3:xix:proposition:3-8}\textbf{\textit{\SGARefLink{sga3:xix:proposition:3:8}{Proposition 3.8}.}}---
\textit{Let \(S\) be a scheme, let \(G\) be an \(S\)-group, and let
\(T\) be a maximal torus of \(G\).  The functor \(\mathcal R\) of
roots of \(G\) with respect to \(T\) is represented by a finite
twisted constant \(S\)-scheme \textup{(\SGARefLink{sga3:x:definition:5-1}{Exposé X, 5.1})} which is an
open and closed subscheme of}
\[
  \underline{\operatorname{Hom}}_{S\text{-gr.}}
  (T,\mathbb G_{m,S}).
\]

\textit{A subset
\(R\subset
\operatorname{Hom}_{S\text{-gr.}}(T,\mathbb G_{m,S})\)
is a root system of \(G\) with respect to \(T\) if and only if the
morphism}
\[
  R_S
  \longrightarrow
  \underline{\operatorname{Hom}}_{S\text{-gr.}}
  (T,\mathbb G_{m,S}),
\]
\textit{defined by the preceding inclusion, induces an isomorphism}
\[
  R_S\xrightarrow{\ \sim\ }\mathcal R.
\]
\SGARefTarget{sga3:xix:proposition:3:8}

\par\medskip\noindent
\textbf{\SGAInlineRefTarget{sga3:xix:paragraph:3-9}3.9.}---
Let \(S\) be a scheme, let \(G\) be a reductive \(S\)-group, let \(T\)
be a maximal torus of \(G\), and let \(\alpha\) be a root of \(G\)
with respect to \(T\), that is, a section of \(\mathcal R\).  Consider
the kernel \(\operatorname{Ker}(\alpha)\) of \(\alpha\), its unique
maximal torus \(T_\alpha\), and the centralizer of the latter,
\[
  Z_\alpha=\underline{\operatorname{Centr}}_G(T_\alpha).
\]
This is a closed \(S\)-subgroup of \(G\), reductive (\SGARefLink{sga3:xix:paragraph:2:8}{2.8}), and of
semisimple rank one (\SGARefLink{sga3:xix:proposition:1-12}{1.12}).  Moreover,
\[
  \operatorname{Lie}(Z_\alpha/S)
  =
  \mathfrak t
  \oplus\mathfrak g^\alpha
  \oplus\mathfrak g^{-\alpha};
\]
hence \(\{\alpha,-\alpha\}\) is a root system of \(Z_\alpha\) with
respect to \(T\).
\SGARefTarget{sga3:xix:paragraph:3:9}

\par\medskip\noindent
\textbf{3.10.}---
Let \(S\) be a scheme, let \(G\) be a reductive \(S\)-group, let \(T\)
be a maximal torus of \(G\), and let \(\alpha\) be a root of \(G\)
with respect to \(T\).  If \(q\in\mathbb Q\) and \(q\alpha\) is a root
of \(G\) with respect to \(T\), then \(q=1\) or \(q=-1\).  This follows
at once from \SGARefLink{sga3:xix:proposition:1-12}{1.12}.
\SGARefTarget{sga3:xix:paragraph:3:10}

% END inputs/XIX/tex/components/14_expose_XIX_proposition3_8_through_3_10_page015.tex

% BEGIN inputs/XIX/tex/components/15_expose_XIX_section4_1_and_lemma4_2_pages015_016.tex
\SGASection{4}{Roots and vector group schemes}
\SGARefTarget{sga3:xix:section:4}

\par\medskip\noindent
\textbf{\SGAInlineRefTarget{sga3:xix:paragraph:4-1}4.1.}---
Let \(S\) be a scheme and let \(F\) be a finite locally free
\(\mathcal O_S\)-module.  The \(S\)-scheme \(\mathbb W(F)\) is smooth
over \(S\).  Its Lie algebra is canonically isomorphic to \(F\).
Indeed, there is a canonical isomorphism
\[
  \mathbb W(F)
  \xrightarrow{\ \sim\ }
  \operatorname{Lie}(\mathbb W(F)/S)
  =
  \mathbb W\bigl(\operatorname{Lie}(\mathbb W(F)/S)\bigr)
\]
\textup{(\SGARefLink{sga3:ii:corollary:4-4-1}{Exposé II, 4.4.1} and \SGARefLink{sga3:xix:scholium:4-4-2}{4.4.2})}.  We shall always identify
\(F\) with \(\operatorname{Lie}(\mathbb W(F)/S)\).
\SGARefTarget{sga3:xix:paragraph:4:1}

\par\medskip\noindent
\textbf{\textit{\SGARefLink{sga3:xix:lemma:4:2}{Lemma 4.2}.}}---
\textit{Let \(S\) be a scheme and let \(V\) be a vector bundle over
\(S\), smooth over \(S\).  There is then a unique isomorphism of
\(\mathcal O_S\)-modules}
\[
  \exp\colon
  \mathbb W\bigl(\operatorname{Lie}(V/S)\bigr)
  \xrightarrow{\ \sim\ } V
\]
\textit{which induces the identity on Lie algebras.}
\footnote{Editor's note: the proof in the original edition has been
retained; it may also be expanded as follows.

Let \(\mathscr F\) be a quasi-coherent \(\mathcal O_S\)-module and put
\(V=\mathbb V(\mathscr F)\).  Write \(\pi\colon V\to S\) for the
projection and \(\epsilon\colon S\to V\) for the zero section.  Then
\(\Omega^1_{V/S}=\pi^*\mathscr F\), whence
\[
  \tag{1}\label{sga3:xix:paragraph:4-1:equation:1}
  \omega^1_{V/S}
  =
  \epsilon^*\Omega^1_{V/S}
  \simeq
  \epsilon^*\pi^*\mathscr F
  \simeq
  \mathscr F,
\]
and therefore
\(\operatorname{Lie}(V/S)=(\omega^1_{V/S})^\vee
\simeq\mathscr F^\vee\).
If \(V\) is smooth over \(S\), then by \textup{(1)}
\(\mathscr F\) is finite locally free, and hence
\[
  \tag{2}\label{sga3:xix:paragraph:4-1:equation:2}
  V=\mathbb V(\mathscr F)
  \simeq
  \mathbb W(\mathscr F^\vee)
  \simeq
  \mathbb W\bigl(\operatorname{Lie}(V/S)\bigr).
\]
Now let \(U\) be an \(S\)-scheme equipped with a section
\(\tau\colon S\to U\), and suppose that an isomorphism
\(\phi\colon V\xrightarrow{\sim}U\) of pointed \(S\)-schemes has been
given, that is, \(\phi\circ\epsilon=\tau\) \textup{(for example, when
\(U\) is an \(S\)-group)}.  Then \(\phi\) induces an isomorphism
\[
  d\phi\colon
  \mathscr F^\vee
  \xrightarrow{\ \sim\ }
  \operatorname{Lie}(U/S),
\]
and therefore an isomorphism \(\psi\) displayed by
\[
\begin{tikzcd}[
  ampersand replacement=\&,
  column sep=huge,
  row sep=large
]
  \mathbb W(\mathscr F^\vee)
    \arrow[r,equal]
  \&
  \mathbb V(\mathscr F)
    \arrow[r,"\phi","\sim"']
  \&
  U
  \\
  \&
  \mathbb W\bigl(\operatorname{Lie}(U/S)\bigr)
    \arrow[ul,"\sim","\mathbb W(d\phi^{-1})"']
    \arrow[ur,"\sim","\psi"']
  \&
\end{tikzcd}
\]
which identifies \(U\) with
\(\mathbb W(\operatorname{Lie}(U/S))\).  Since the functor
\(\mathbb W\) is fully faithful, there is a unique isomorphism
\[
  \exp\colon
  \mathbb W\bigl(\operatorname{Lie}(U/S)\bigr)
  \xrightarrow{\ \sim\ }U
\]
such that
\(\operatorname{Lie}(\psi^{-1}\circ\exp)=\operatorname{id}\).
In fact, one sees directly that
\(\operatorname{Lie}(\psi)=\operatorname{id}\), so that
\(\exp=\psi\).}
\SGARefTarget{sga3:xix:lemma:4:2}

\SGAProof{Proof}
We have \(V=\mathbb V(\mathscr F)\) for some quasi-coherent
\(\mathcal O_S\)-module \(\mathscr F\).  Since \(V\) is smooth over
\(S\), the module
\(\mathscr F\simeq\omega^1_{V/S}\) is finite locally free, and hence
\[
  \operatorname{Lie}(V/S)
  =
  \mathbb V(\omega^1_{V/S})
  \simeq
  \mathbb W\bigl(\operatorname{Lie}(V/S)\bigr).
\]
Moreover, by Exposé II, loc.\ cit., there is a canonical isomorphism
\[
  V
  \xrightarrow{\ \sim\ }
  \operatorname{Lie}(V/S)
  \simeq
  \mathbb W\bigl(\operatorname{Lie}(V/S)\bigr),
\]
and the uniqueness of \(\exp\) follows at once because
\(\mathbb W\) is fully faithful.
\hfill\(\square\)

% END inputs/XIX/tex/components/15_expose_XIX_section4_1_and_lemma4_2_pages015_016.tex

% BEGIN inputs/XIX/tex/components/16_expose_XIX_notations4_3_through_scholium4_4_2_page016.tex
\par\medskip\noindent
\textbf{\SGAInlineRefTarget{sga3:xix:paragraph:4-3}4.3. Notation.}---
If \(V\) is a vector bundle over \(S\), let \(V^\times\) denote the
open subset obtained from \(V\) by removing the zero section.  Write
the group law of \(V\) multiplicatively.  The action of
\(\mathcal O_S\) on \(V\) which defines its module structure will
then be written exponentially:
\[
  \mathcal O_S\times_S V\longrightarrow V,
  \qquad
  (x,v)\longmapsto v^x.
\]
Thus
\[
  (vv')^x=v^x(v')^x,
  \qquad
  v^{x+x'}=v^xv^{x'},
  \qquad
  v^{xx'}=(v^x)^{x'}.
\]
In particular, upon restricting the operator action to
\(\mathbb G_{m,S}\), the open subset \(V^\times\) is stable and is
therefore an object with operator group \(\mathbb G_{m,S}\).  We
shall continue to write this action as
\((z,v)\longmapsto v^z\).
\SGARefTarget{sga3:xix:paragraph:4:3}

\par\medskip\noindent
\SGAInlineRefTarget{sga3:xix:definition:4-4-1}\textbf{\textit{\SGARefLink{sga3:xix:definition:4:4:1}{Definition 4.4.1}.}}\footnote{Editor's note: for
purposes of later references, paragraph 4.4 has been divided into
\SGARefLink{sga3:xix:definition:4-4-1}{4.4.1} and \SGARefLink{sga3:xix:scholium:4-4-2}{4.4.2}.}---
\textit{Let \(L\) be an invertible module on \(S\), and let
\(\mathbb W(L)\) be the associated vector bundle.  Then
\(\mathbb W(L)^\times\) is a \textup{(locally trivial)} principal
homogeneous bundle under \(\mathbb G_{m,S}\).  We write}
\[
  \Gamma(S,L)^\times=\mathbb W(L)^\times(S).
\]
\SGARefTarget{sga3:xix:definition:4:4:1}

\par\medskip\noindent
\SGAInlineRefTarget{sga3:xix:scholium:4-4-2}\textbf{\textit{\SGARefLink{sga3:xix:scholium:4:4:2}{Scholium 4.4.2}.}}---
\textit{There are natural bijections among the isomorphisms of
\(\mathcal O_S\)-modules
\(\mathcal O_S\simeq L\), the isomorphisms of
\(\mathcal O_S\)-modules
\(\mathcal O_S\simeq\mathbb W(L)\),\footnote{Editor's note: here the
vector bundle \(\mathbb W(L)\) is identified with the
\(\mathcal O_S\)-module functor which it represents.} and the
sections \(S\to\mathbb W(L)^\times\).}

The correspondence is
\[
  f\longmapsto\mathbb W(f)\longmapsto\mathbb W(f)(1).
\]
It is compatible with base change.  Thus
\(\mathbb W(L)^\times\) may be regarded as the ``scheme of
trivializations of \(\mathbb W(L)\)''.
\SGARefTarget{sga3:xix:scholium:4:4:2}


% END inputs/XIX/tex/components/16_expose_XIX_notations4_3_through_scholium4_4_2_page016.tex

% BEGIN inputs/XIX/tex/components/17_expose_XIX_paragraph4_5_through_corollary4_9_page017.tex
\par\medskip\noindent
\textbf{\SGAInlineRefTarget{sga3:xix:paragraph:4-5}4.5.}---
Let \(S\) be a scheme, let \(T\) be an \(S\)-torus, let \(P\) be an
\(S\)-group with operator group \(\mathbb G_{m,S}\)
\textup{(for example, a vector bundle)}, and let \(\alpha\) be a
character of \(T\).  We write \(T\mathbin{\cdot_\alpha}P\) for the
semidirect product of \(P\) by \(T\), where \(T\) acts on \(P\)
through the composite morphism
\[
  T
  \xrightarrow{\ \alpha\ }
  \mathbb G_{m,S}
  \longrightarrow
  \underline{\operatorname{Aut}}_{S\text{-gr.}}(P).
\]
\SGARefTarget{sga3:xix:paragraph:4:5}

\par\medskip\noindent
\SGAInlineRefTarget{sga3:xix:definition:4-6}\textbf{\textit{\SGARefLink{sga3:xix:definition:4:6}{Definition 4.6}.}}---
\textit{Let \(S\) be a scheme, let \(G\) be an \(S\)-group scheme,
let \(T\) be a subgroup of \(G\), let \(\alpha\) be a character of
\(T\), and let \(L\) be an \(\mathcal O_S\)-module.  Let}
\[
  p\colon\mathbb W(L)\longrightarrow G
\]
\textit{be a morphism of \(S\)-group functors.}\footnote{Editor's
note: the edition uses a bold \(\mathbf W(L)\) here because, for an
arbitrary \(\mathcal O_S\)-module \(L\), the functor need not be
representable.  Below, however, \(L\) is assumed finite locally free
\textup{(indeed, invertible)}, in which case the functor is
represented by the vector bundle \(\mathbb V(L^\vee)\), and the
notation \(\mathbb W(L)\) is used without this distinction.}
\textit{We say that \(p\) is \emph{normalized by \(T\) with
multiplier \(\alpha\)} if it satisfies the following equivalent
conditions:}

\begin{enumerate}
\item[\textup{(i)}]
\(p\) is a morphism of objects with operator group \(T\), where \(T\)
acts on \(\mathbb W(L)\) through \(\alpha\), and on \(G\) by inner
automorphisms.  In other words, for every \(S'\to S\), every
\(t\in T(S')\), and every
\[
  x\in\mathbb W(L)(S')
  =
  \Gamma(S',L\otimes_{\mathcal O_S}\mathcal O_{S'}),
\]
one has
\[
  \operatorname{int}(t)p(x)=p\bigl(\alpha(t)x\bigr).
\]

\item[\textup{(ii)}]
The morphism
\[
  T\mathbin{\cdot_\alpha}\mathbb W(L)\longrightarrow G,
  \qquad
  (t,x)\longmapsto t\,p(x),
\]
defined by multiplication in \(G\), is a group morphism.
\end{enumerate}
\SGARefTarget{sga3:xix:definition:4:6}

\par\medskip\noindent
\SGAInlineRefTarget{sga3:xix:lemma:4-7}\textbf{\textit{\SGARefLink{sga3:xix:lemma:4:7}{Lemma 4.7}.}}---
\textit{Under the hypotheses of \SGARefLink{sga3:xix:definition:4-6}{4.6}, suppose in addition that \(G\)
is smooth with connected fibers,\footnote{Editor's note: compare
editor's note \textup{(28)}.} that \(T\) is a maximal torus of \(G\),
and that \(L\) is invertible.  If \(p\) is a monomorphism and
\(\alpha\) is nonzero on every fiber, then \(\alpha\) is a root of
\(G\) with respect to \(T\).}
\SGARefTarget{sga3:xix:lemma:4:7}

Indeed,
\[
  \operatorname{Lie}(p)\colon L\longrightarrow\mathfrak g
\]
is a monomorphism which factors through \(\mathfrak g^\alpha\).

\par\medskip\noindent
\SGAInlineRefTarget{sga3:xix:proposition:4-8}\textbf{\textit{\SGARefLink{sga3:xix:proposition:4:8}{Proposition 4.8}.}}---
\textit{Under the hypotheses of \SGARefLink{sga3:xix:lemma:4-7}{4.7}, suppose that \(G\) is reductive
and that \(p\) is a monomorphism.  Then \(\alpha\) is a root of \(G\)
with respect to \(T\), and \(\operatorname{Lie}(p)\) induces an
isomorphism}
\[
  \operatorname{Lie}(p)\colon
  L\xrightarrow{\ \sim\ }\mathfrak g^\alpha.
\]
\SGARefTarget{sga3:xix:proposition:4:8}

Indeed, by \SGARefLink{sga3:xix:lemma:4-7}{4.7} and the fact that \(\mathfrak g^\alpha\) is
invertible, it is enough to prove that \(\alpha\) is nonzero on every
fiber.  Let \(s\in S\) be such that
\(\alpha_{\bar s}=0\) \textup{(\(=1\) in multiplicative notation)}.
If \(X\) is a nonzero section of
\(L\otimes_S\kappa(\bar s)\), then \(p(X)\) is a unipotent element
different from the identity in \(G(\bar s)\) which centralizes
\(T_{\bar s}\).  This is impossible, since \(T_{\bar s}\) is its own
centralizer.

\par\medskip\noindent
\SGAInlineRefTarget{sga3:xix:corollary:4-9}\textbf{\textit{\SGARefLink{sga3:xix:corollary:4:9}{Corollary 4.9}.}}---
\textit{Under the hypotheses of \SGARefLink{sga3:xix:proposition:4-8}{4.8}, there exists a monomorphism of
groups with operator group \(T\) \textup{(that is, normalized by
\(T\) with multiplier \(\alpha\))}
\[
  \mathbb W(\mathfrak g^\alpha)\longrightarrow G
\]
which induces on Lie algebras the canonical morphism
\(\mathfrak g^\alpha\to\mathfrak g\).}
\SGARefTarget{sga3:xix:corollary:4:9}

We shall soon see that \SGARefLink{sga3:xix:corollary:4-9}{4.9} in fact holds whenever \(G\) is a
reductive group and \(\alpha\) is a root of \(G\) with respect to
\(T\), and that such a morphism is unique.


% END inputs/XIX/tex/components/17_expose_XIX_paragraph4_5_through_corollary4_9_page017.tex

% BEGIN inputs/XIX/tex/components/18_expose_XIX_reminder4_10_through_proposition4_12_statement_page018.tex
\par\medskip\noindent
\textbf{\textit{Reminder 4.10.}}---
\textit{Let \(k\) be an algebraically closed field, let \(G\) be a
reductive \(k\)-group, let \(T\) be a maximal torus of \(G\), and let
\(\alpha\) be a root of \(G\) with respect to \(T\).  There exists a
monomorphism}
\[
  p\colon\mathbb G_{a,k}\longrightarrow G
\]
\textit{normalized by \(T\) with multiplier \(\alpha\).}

See the \emph{Bible}, §13.1, Theorem 1.
\SGARefTarget{sga3:xix:reminder:4:10}

\par\medskip\noindent
\textbf{\SGAInlineRefTarget{sga3:xix:paragraph:4-11}4.11.}---
We finish this section with a technical result which will be useful
to us.  Let \(S\) be a scheme, and let \(L\) be an invertible
\(\mathcal O_S\)-module.  Let \(q>0\) be an integer such that
\(x\longmapsto x^q\) is an endomorphism of the \(S\)-group
\(\mathbb G_{a,S}\).  If \(S\ne\varnothing\), then \(q=1\), or
\(q=p^n\), where \(p\) is a prime number which vanishes on \(S\).
This follows at once from the elementary fact that the greatest
common divisor of the binomial coefficients
\(\binom qi\), for \(i\ne0,q\), is \(p\) if \(q=p^n\) with \(p\)
prime, and is \(1\) otherwise.

The morphism defined by taking the \(q\)-th power,
\[
  L\longrightarrow L^{\otimes q},
\]
is a morphism of sheaves of abelian groups.  By base change it
defines a morphism of \(S\)-group schemes
\[
  \mathbb W(L)\longrightarrow\mathbb W(L^{\otimes q}).
\]
In particular, if \(L'\) is another invertible module and an
\(\mathcal O_S\)-linear morphism
\[
  h\colon L^{\otimes q}\longrightarrow L'
\]
is given, one obtains a morphism of \(S\)-group schemes
\[
  \mathbb W(L)\longrightarrow\mathbb W(L'),
  \qquad
  x\longmapsto h(x^q).
\]
\SGARefTarget{sga3:xix:paragraph:4:11}

With this notation, we have the following result.

\par\medskip\noindent
\textbf{\textit{\SGARefLink{sga3:xix:proposition:4:12}{Proposition 4.12}.}}---
\textit{Let \(S\) be a scheme, let \(T\) \textup{(resp. \(T'\))} be
an \(S\)-torus, let \(L\) \textup{(resp. \(L'\))} be an invertible
\(\mathcal O_S\)-module, and let \(\alpha\)
\textup{(resp. \(\alpha'\))} be a character of \(T\)
\textup{(resp. \(T'\))}.  Let \(f\colon T\to T'\) be a group
morphism and let
\[
  g\colon\mathbb W(L)\longrightarrow\mathbb W(L')
\]
be a morphism of \(S\)-schemes \textup{(not necessarily a morphism
of \(S\)-groups)} satisfying}
\[
  \tag{\(\star\)}\label{sga3:xix:paragraph:4-11:equation:star}
  g\bigl(\alpha(t)x\bigr)
  =
  \alpha'\bigl(f(t)\bigr)g(x)
\]
\textit{for every \(S'\to S\), every \(x\in\mathbb W(L)(S')\), and
every \(t\in T(S')\).  Let \(s_0\in S\) be such that
\(\alpha_{\bar s_0}\ne0\).}
\footnote{Editor's note: the phrases ``consider the semidirect
product \(T\mathbin{\cdot_\alpha}\mathbb W(L)\), respectively
\(T'\mathbin{\cdot_{\alpha'}}\mathbb W(L')\)'' have been deleted.
On the other hand, in part \textup{\SGARefLink{sga3:xix:paragraph:4:11:item:b}{(b)}} the hypothesis that \(g\) be
a group morphism, together with \textup{\((\star)\)}, is equivalent
to requiring
\[
  (t,x)\longmapsto\bigl(f(t),g(x)\bigr)
\]
to be a group morphism from
\(T\mathbin{\cdot_\alpha}\mathbb W(L)\) to
\(T'\mathbin{\cdot_{\alpha'}}\mathbb W(L')\).}

\begin{enumerate}
\item[\textup{(a)}]\SGAInlineRefTarget{sga3:xix:proposition:4:12:item:a}
\textit{Suppose that \(g\) sends the zero section to the zero section
and that, for every integer \(n>0\),}
\[
  (\alpha'\circ f)_{\bar s_0}\ne n\alpha_{\bar s_0}.
\]
\textit{Then \(g=0\) in a neighborhood of \(s_0\).}

\item[\textup{\SGAInlineRefTarget{sga3:xix:proposition:4-12:item:b}\SGAInlineRefTarget{sga3:xix:paragraph:4:11:item:b}(b)}]
\textit{Suppose that \(g\) is a group morphism and that
\(g_{\bar s_0}\ne0\).  There then exist an open subset \(U\subset S\)
containing \(s_0\) and an integer \(q>0\) such that
\(x\longmapsto x^q\) is an endomorphism of
\(\mathbb G_{a,U}\) and}
\[
  (\alpha'\circ f)_U=q\alpha_U.
\]
\end{enumerate}
\SGARefTarget{sga3:xix:proposition:4:12}


% END inputs/XIX/tex/components/18_expose_XIX_reminder4_10_through_proposition4_12_statement_page018.tex

% BEGIN inputs/XIX/tex/components/19_expose_XIX_proposition4_12_part_c_and_proof_page019.tex
\begin{enumerate}
\item[\textup{\SGAInlineRefTarget{sga3:xix:proposition:4:12:item:c}(c)}]
\textit{Suppose that}
\[
  (\alpha'\circ f)_{\bar s_0}=q\alpha_{\bar s_0},
\]
\textit{where \(q>0\) is an integer such that
\(x\longmapsto x^q\) is an endomorphism of
\(\mathbb G_{a,S}\).  There then exist an open subset \(U\subset S\)
containing \(s_0\) and a unique morphism of
\(\mathcal O_S\)-modules}
\[
  h\colon L^{\otimes q}|_U\longrightarrow L'|_U
\]
\textit{such that \(g_U\) is the composite}
\[
  \mathbb W(L)_U
  \xrightarrow{\ x\mapsto x^q\ }
  \mathbb W(L^{\otimes q})_U
  \xrightarrow{\ \mathbb W(h)\ }
  \mathbb W(L')_U.
\]
\end{enumerate}

\SGAProof{Proof of \SGARefLink{sga3:xix:proposition:4:12}{Proposition 4.12}.}
Let us prove \textup{(a)}.  Since the conclusion is local on \(S\),
we may suppose that
\(\mathbb W(L)=\mathbb W(L')=\mathbb G_{a,S}\), and hence that \(g\)
is given by a polynomial
\[
  g(X)=\sum_{n\geq0}a_nX^n,
  \qquad
  a_n\in\Gamma(S,\mathcal O_S).
\]
Condition \textup{\((\star)\)}, which relates \(f\) and \(g\), is the
identity in \(\Gamma(S',\mathcal O_{S'})[X]\)
\[
  \sum_{n\geq0}a_n\alpha'\bigl(f(t)\bigr)X^n
  =
  \sum_{n\geq0}a_n\alpha(t)^nX^n.
\]
Thus, for every \(n\geq0\), every \(S'\to S\), and every
\(t\in T(S')\),
\[
  a_n\bigl(\alpha'(f(t))-\alpha(t)^n\bigr)=0.
\]

For every \(n\geq0\), let \(S_n\) be the set of points \(s\in S\)
such that
\[
  (\alpha'\circ f)_{\bar s}=n\alpha_{\bar s}.
\]
It is known \textup{(\SGARefLink{sga3:ix:corollary:5:3}{Exposé IX, 5.3})} that the \(S_n\) are both open
and closed and that
\[
  (\alpha'\circ f)_{S_n}=n\alpha_{S_n}.
\]
Moreover, since \(\alpha_{\bar s_0}\ne0\), after shrinking \(S\) we
may suppose that \(\alpha\) is nonzero on every fiber \textup{(same
reference)}; it follows that the \(S_n\) are pairwise disjoint.
After shrinking \(S\) once more, we may therefore suppose that one
of the following two cases holds: \(S=S_n\) for some \(n\), or every
\(S_n\) is empty.

Let \(m\geq0\) be such that \(S_m=\varnothing\).  We claim that
\(a_m=0\).  Indeed, \(\alpha'\circ f\) and \(m\alpha\) are distinct
on every fiber of \(S\), and we have:

\par\medskip\noindent
\textbf{\textit{\SGARefLink{sga3:xix:lemma:4:13}{Lemma 4.13}.}}---
\textit{Let \(S\) be a scheme, let \(T\) be an \(S\)-torus, and let
\(\alpha,\alpha'\) be two characters of \(T\) which are distinct on
every fiber.  There exist an fpqc covering family
\(\{S_i\to S\}\) and, for every \(i\), an element \(t_i\in T(S_i)\)
such that}
\[
  \alpha(t_i)-\alpha'(t_i)=1.
\]
\SGARefTarget{sga3:xix:lemma:4:13}

The lemma is immediate by reduction first to the diagonalizable
case, and then to \(T=\mathbb G_{m,S}\).

We return to the proof of \SGARefLink{sga3:xix:proposition:4:12}{Proposition 4.12}.  We have just proved
\textup{(a)}.  In cases \textup{(b)} and \textup{\SGARefLink{sga3:xix:proposition:4:12:item:c}{(c)}}, there is an
integer \(n\) such that \(S=S_n\) \textup{(\(n=q\) in
\textup{\SGARefLink{sga3:xix:proposition:4:12:item:c}{(c)}})}.  The preceding argument gives \(a_m=0\) for
\(m\ne n\), and therefore
\[
  g(X)=a_nX^n,
  \qquad
  a_n\in\Gamma(S,\mathcal O_S).
\]
This proves \textup{(c)} at once.  In case \textup{(b)}, we know
that \(a_n(s_0)\ne0\); we may therefore suppose that \(a_n\) is
invertible on \(S\).  It follows that \(x\longmapsto x^n\) is an
endomorphism of \(\mathbb G_{a,S}\), by the hypothesis in
\textup{(b)}, and this completes the proof.


% END inputs/XIX/tex/components/19_expose_XIX_proposition4_12_part_c_and_proof_page019.tex

% BEGIN inputs/XIX/tex/components/20_expose_XIX_section5_1_through_lemma5_3_page020.tex
\SGANextSectionTarget{sga3:xix:section:5}\SGASection{5}{An instructive example}

\par\medskip\noindent
\textbf{\SGAInlineRefTarget{sga3:xix:paragraph:5-1}5.1.}---
Let \(k\) be an algebraically closed field of characteristic \(0\).
Put \(A=k[t]\), the polynomial ring in one variable over \(k\), and
let \(S=\operatorname{Spec}(A)\).  Consider the following Lie algebra
\(\mathfrak g\) over \(\mathcal O_S\).  As a module it is free of
rank \(3\), with basis \(\{X,Y,H\}\), and its multiplication table is
\[
  [X,Y]=2tH,
  \qquad
  [H,X]=X,
  \qquad
  [H,Y]=-Y.
\]
For \(s\in S\), \(s\ne s_0\), where \(s_0\) is the point defined by
\(t=0\), the fiber
\[
  \mathfrak g(s)=\mathfrak g\otimes_A\kappa(s)
\]
is isomorphic to the Lie algebra of
\(\mathrm{PGL}_{2,\kappa(s)}\).  For \(s=s_0\), it is a solvable Lie
algebra.
\SGARefTarget{sga3:xix:paragraph:5:1}

\par\medskip\noindent
\textbf{\SGAInlineRefTarget{sga3:xix:paragraph:5-2}5.2.}---
Let \(G_1\) be the group scheme of automorphisms of \(\mathfrak g\):
for every \(S'\to S\), \(G_1(S')\) is the group of automorphisms of
the \(\mathcal O_{S'}\)-Lie algebra
\(\mathfrak g\otimes_{\mathcal O_S}\mathcal O_{S'}\).  It is a closed
subscheme of the group \(\mathrm{GL}(\mathfrak g)\) of automorphisms
of the \(\mathcal O_S\)-module \(\mathfrak g\).

Let \(S'\to S\), and let
\[
  u\in M_3\bigl(\Gamma(S',\mathcal O_{S'})\bigr)
\]
be viewed as an endomorphism of the \(\mathcal O_{S'}\)-module
\(\mathfrak g\otimes_{\mathcal O_S}\mathcal O_{S'}\):
\[
  \begin{aligned}
  u(X)&=aX+bY+eH,\\
  u(Y)&=b'X+a'Y+e'H,\\
  u(H)&=cX+c'Y+dH.
  \end{aligned}
\]
It is immediate that \(u\) is a section of \(G_1\) if and only if
\(\det(u)\) is invertible and the following relations hold:
\[
\begin{alignedat}{4}
  \textup{(1)}\quad&a(d-1)=ec,
  &\qquad&
  \textup{(1')}\quad&a'(d-1)=e'c',\\
  \textup{(2)}\quad&b(d+1)=ec',
  &&
  \textup{(2')}\quad&b'(d+1)=e'c,\\
  \textup{(3)}\quad&e=2t(bc-ac'),
  &&
  \textup{(3')}\quad&e'=2t(b'c'-a'c),\\
  \textup{(4)}\quad&2tc=eb'-ae',
  &&
  \textup{(4')}\quad&2tc'=be'-ea',\\
  \textup{(5)}\quad&2t(aa'-bb')=2td.&&&
\end{alignedat}
\]
\footnote{Editor's note: the equality
\([u(X),u(Y)]=2t\,u(H)\) gives relations \textup{(4)},
\textup{(4')}, and \textup{(5)}; the equality
\([u(H),u(X)]=u(X)\) gives \textup{(1)--(3)}; and the equality
\([u(H),u(Y)]=-u(Y)\) gives \textup{(1')--(3')}.}
\SGARefTarget{sga3:xix:paragraph:5:2}

\par\medskip\noindent
\SGAInlineRefTarget{sga3:xix:lemma:5-3}\textbf{\textit{\SGARefLink{sga3:xix:lemma:5:3}{Lemma 5.3}.}}---
\textit{Relations \textup{(1)}, \textup{(1')}, \textup{(2)}, and
\textup{(2')} imply}
\[
  \det(u)=aa'(2-d)+bb'(d+2),
  \qquad
  aa'-bb'=d\det(u).
\]
\SGARefTarget{sga3:xix:lemma:5:3}

Indeed, the first assertion follows at once by substituting
relations \textup{(1)}, \textup{(1')}, \textup{(2)}, and
\textup{(2')} in the expansion of \(\det(u)\):
\[
\begin{aligned}
  \det(u)
  &=aa'd+be'c+b'c'e-a'ec-ae'c'-bb'd\\
  &=aa'd+bb'(d+1)+bb'(d+1)
    -aa'(d-1)-bb'd-aa'(d-1)\\
  &=aa'(d-d+1-d+1)+bb'(d+1+d+1-d)\\
  &=aa'(2-d)+bb'(d+2).
\end{aligned}
\]
Multiplying this relation by \(d\), we obtain
\[
  d\det(u)=aa'(2d-d^2)+bb'(d^2+2d).
\]


% END inputs/XIX/tex/components/20_expose_XIX_section5_1_through_lemma5_3_page020.tex

% BEGIN inputs/XIX/tex/components/21_expose_XIX_lemma5_3_completion_through_lemma5_7_page021.tex
On the other hand, the relation
\(\textup{(1)}\times\textup{(1')}=
\textup{(2)}\times\textup{(2')}\) gives at once
\[
  aa'(d-1)^2=bb'(d+1)^2.
\]
Combining the last two relations immediately yields the second
formula of \SGARefLink{sga3:xix:lemma:5:3}{Lemma 5.3}.

\par\medskip\noindent
\textbf{5.4.}---
Now consider
\[
  G=G_1\cap\mathrm{SL}(\mathfrak g).
\]
This is the closed subgroup of \(G_1\) defined by the equation
\(\det(u)=1\); hence it is an affine group over \(S\).
\SGARefTarget{sga3:xix:paragraph:5:4}

\par\medskip\noindent
\SGAInlineRefTarget{sga3:xix:proposition:5-5}\textbf{\textit{\SGARefLink{sga3:xix:proposition:5:5}{Proposition 5.5}.}}---
\textit{The group \(G\) is smooth over \(S\).}
\SGARefTarget{sga3:xix:proposition:5:5}

To prove the proposition, we shall need the following lemmas.

\par\medskip\noindent
\textbf{\textit{\SGARefLink{sga3:xix:lemma:5:6}{Lemma 5.6}.}}---
\textit{Let \(U\) be the open subset of}
\[
  \operatorname{End}_A(\mathfrak g)
  \simeq
  \mathbb W\bigl(M_3(\mathcal O_S)\bigr)
\]
\textit{defined by the condition that \(ad\) be invertible; that is,
the principal open subset
\(\operatorname{End}_A(\mathfrak g)_f\), where
\(f(u)=ad\).  Then \(U\cap G\) is the closed subscheme of \(U\)
defined by the six equations}
\[
  \textup{(1)},\quad
  \textup{(2)},\quad
  \textup{(2')},\quad
  \textup{(3)},\quad
  \textup{(3')},\quad
  \textup{(D)}\colon aa'-bb'=d.
\]
\SGARefTarget{sga3:xix:lemma:5:6}

It is clear that these six relations are satisfied by every point
of \(G\), by \SGARefLink{sga3:xix:lemma:5-3}{Lemma 5.3}, and therefore by every point of \(U\cap G\).
Conversely, suppose that \(u\in U(S')\), for some \(S'\to S\),
satisfies the six conditions in the statement.  We must show that
\(\det(u)=1\), and that \(u\) also satisfies \textup{(1')},
\textup{(4)}, \textup{(4')}, and \textup{(5)}.

First, \textup{(D)} implies \textup{(5)}.  Relations \textup{(2)}
and \textup{(2')} give
\[
  bb'(d+1)=bce'=b'c'e.
\]
On the other hand, from \textup{(3)} and \textup{(3')}, writing
\(2t(bc-ac')(b'c'-a'c)\) in two ways, we obtain
\[
  (bc-ac')e'=(b'c'-a'c)e.
\]
Together with the preceding relation, this gives
\[
  ac'e'=a'ce.
\]
By \textup{(1)}, however,
\[
  a'ce=a'a(d-1).
\]
Thus
\[
  a\bigl(a'(d-1)-e'c'\bigr)=0,
\]
and this implies \textup{(1')}, since \(a\) is invertible.

Consequently, \textup{(1)}, \textup{(2)}, \textup{(2')}, and
\textup{(1')} hold.  \SGARefLink{sga3:xix:lemma:5:3}{Lemma 5.3} and \textup{(D)} now give
\[
  d\bigl(\det(u)-1\bigr)=0.
\]
Since \(d\) is invertible, it follows that \(\det(u)=1\).

It remains to prove \textup{(4)} and \textup{(4')}.  We prove
\textup{(4')}; the other calculation follows by symmetry.  From
\textup{(3)}, \textup{(3')}, and \textup{(D)}, we have
\[
  a'e+be'=-2t(aa'-bb')c'=-2tdc'.
\]
Combining \textup{(1')} and \textup{(2)}, we also have
\[
  a'e+be'=-d(be'-ea'),
\]
\footnote{Editor's note: the sign has been corrected.}
which proves \textup{(4')}, since \(d\) is invertible.

\par\medskip\noindent
\SGAInlineRefTarget{sga3:xix:lemma:5-7}\textbf{\textit{\SGARefLink{sga3:xix:lemma:5:7}{Lemma 5.7}.}}---
\textit{The group \(G\) is smooth over \(S\) along the identity
section.}
\SGARefTarget{sga3:xix:lemma:5:7}


% END inputs/XIX/tex/components/21_expose_XIX_lemma5_3_completion_through_lemma5_7_page021.tex

% BEGIN inputs/XIX/tex/components/22_expose_XIX_lemma5_7_proof_through_paragraph5_11_page022.tex
\SGAProof{Proof}
By 5.6 and SGA 1, II 4.10, it is enough to prove that the
differentials of the functions
\[
\begin{gathered}
  a(d-1)-ec,\qquad
  b(d+1)-ec',\qquad
  b'(d+1)-e'c,\\
  e-2t(bc-ac'),\qquad
  e'-2t(b'c'-a'c),\qquad
  aa'-bb'-d
\end{gathered}
\]
at the points of the identity section of \(G\) are linearly
independent.  If the differential of a lower-case letter is denoted
by the corresponding capital letter, these differentials are
\[
  D,\qquad
  2B,\qquad
  2B',\qquad
  E+2tC',\qquad
  E'+2tC,\qquad
  A+A'-D,
\]
\footnote{Editor's note: \(A+A'+B\) has been corrected to
\(A+A'-D\).}
and they are indeed linearly independent modulo every
\((t-\lambda)\), \(\lambda\in k\).
\footnote{Editor's note: ``\((t-a)\), \(a\in A\)'' has been
corrected to ``\((t-\lambda)\), \(\lambda\in k\)''.}

\par\medskip\noindent
\textbf{\textit{\SGARefLink{sga3:xix:lemma:5:8}{Lemma 5.8}.}}---
\textit{For \(s\in S\), \(s\ne s_0\), the fiber \(G_s\) is connected
and semisimple.}
\SGARefTarget{sga3:xix:lemma:5:8}

\SGAProof{Proof}
Indeed, since \(s\ne s_0\), \(\mathfrak g(s)\) is isomorphic to the
Lie algebra of \(\mathrm{PGL}_{2,\kappa(s)}\), and, on the other hand,
\(G_s=(G_1)_s\).  The automorphism group of the Lie algebra of
\(\mathrm{PGL}_2\) over a field of characteristic \(0\) is known to
be \(\mathrm{PGL}_2\) itself, which is connected and semisimple.
\footnote{Editor's note: the preceding sentence has been slightly
modified.}

\par\medskip\noindent
\textbf{\textit{\SGARefLink{sga3:xix:lemma:5:9}{Lemma 5.9}.}}---
\textit{The fiber \(G_{s_0}\) is solvable and has two connected
components, of the following form:}
\[
\begin{aligned}
  G^0_{s_0}
  &=
  \left\{
  \begin{pmatrix}
    a&0&0\\
    0&a^{-1}&0\\
    c&c'&1
  \end{pmatrix}
  \right\},
  &
  G^-_{s_0}
  &=
  \left\{
  \begin{pmatrix}
    0&b&0\\
    b^{-1}&0&0\\
    c&c'&-1
  \end{pmatrix}
  \right\}.
\end{aligned}
\]
\SGARefTarget{sga3:xix:lemma:5:9}

\SGAProof{Proof}
We have \(e=e'=0\), since \(t=0\) at \(s_0\).  The equations
\textup{(1)}, \textup{(1')}, \(\ldots\), \textup{(5)}, and
\textup{(D)} can then be solved immediately.

\par\medskip\noindent
\SGAInlineRefTarget{sga3:xix:lemma:5-10}\textbf{\textit{\SGARefLink{sga3:xix:lemma:5:10}{Lemma 5.10}.}}---
\textit{The matrix}
\[
  w=
  \begin{pmatrix}
    0&1&0\\
    1&0&0\\
    0&0&-1
  \end{pmatrix}
\]
\textit{is a section of \(G\) over \(S\), such that
\(w(s_0)\in G^-_{s_0}\).}
\SGARefTarget{sga3:xix:lemma:5:10}

Let us now prove \SGARefLink{sga3:xix:proposition:5-5}{Proposition 5.5}.  Denote by \(G^0\) the union of the
identity components of the fibers of \(G\), that is, the complement
of \(G^-_{s_0}\).  Since \(G\) is smooth over \(S\) along the identity
section by \SGARefLink{sga3:xix:lemma:5-7}{Lemma 5.7}, \(G^0\) is an open subgroup of \(G\), smooth
over \(S\), by Exposé VI\textsubscript{B}, \SGARefLink{thm:sga3-VIB-3.10}{3.10}.  By translation,
\(G\) is also plainly smooth at the points of \(w(S)\); hence \(G\)
is smooth over \(S\).
\footnote{Editor's note: the preceding sentence has been modified,
and the reference to Exposé VI\textsubscript{B}, \SGARefLink{thm:sga3-VIB-3.10}{3.10} has been
added.}

\par\medskip\noindent
\textbf{\SGAInlineRefTarget{sga3:xix:paragraph:5-11}5.11.}---
Consider the morphism
\[
  \mathbb G_{m,S}\longrightarrow G^0,
  \qquad
  z\longmapsto
  \begin{pmatrix}
    z&0&0\\
    0&z^{-1}&0\\
    0&0&1
  \end{pmatrix}.
\]
It is a monomorphism defining a torus \(T\) of \(G^0\).  We claim
that
\[
  T=\operatorname{Centr}_G(T)
   =\operatorname{Centr}_{G^0}(T).
\]
\SGARefTarget{sga3:xix:paragraph:5:11}

% END inputs/XIX/tex/components/22_expose_XIX_lemma5_7_proof_through_paragraph5_11_page022.tex

% BEGIN inputs/XIX/tex/components/23_expose_XIX_paragraph5_11_completion_through_proposition6_1_opening_page023.tex
It is enough to verify the first equality.  Since the two sides are
smooth \(S\)-subgroups of \(G\), it is enough to check that they have
the same geometric points.  In the fibers over points \(s\ne s_0\),
this follows because \(\mathrm{PGL}_{2,\kappa(s)}\) is reductive and
\(T_s\) is a maximal torus in it for dimensional reasons
\textup{(cf.\ \SGARefLink{sga3:xix:lemma:1-11}{1.11})}.  In the fiber over \(s_0\), the verification is
immediate.  In particular, \(T\) is a maximal torus of both \(G\)
and \(G^0\).

\par\medskip\noindent
\textbf{5.12.}---
The section \(w\) of \(G\) defined in \SGARefLink{sga3:xix:lemma:5-10}{Lemma 5.10} normalizes \(T\).
It follows immediately \textup{(cf.\ \SGARefLink{sga3:xix:remark:2-4}{2.4})} that the Weyl group of
\(G\) is isomorphic to \((\mathbb Z/2\mathbb Z)_S\), and is therefore
finite over \(S\):
\[
  W_G(T)=\operatorname{Norm}_G(T)/T
        =(\mathbb Z/2\mathbb Z)_S.
\]
By contrast, \(W_{G^0}(T)\) is not finite over \(S\): it ``lacks one
point'' above \(s_0\).
\SGARefTarget{sga3:xix:paragraph:5:12}

\par\medskip\noindent
\textbf{\SGAInlineRefTarget{sga3:xix:paragraph:5-13}5.13.}---
The open immersion \(G^0\to G\) is not a closed immersion, since
\(G^0\) is dense in \(G\); it is nevertheless an affine morphism
\textup{(and hence \(G^0\) is affine over \(S\))}.  Indeed,
\(G^0_{s_0}\) is closed in \(G_{s_0}\), which is closed in \(G\), so
its complement \(U\) in \(G\) is open.  The open subsets \(G^0\) and
\(U\) cover \(G\), and it is enough to verify that the immersions
\[
  G^0\longrightarrow G^0
  \qquad\text{and}\qquad
  G^0\cap U\longrightarrow U
\]
are affine.  This is trivial for the first; for the second, observe
that \(U\cap G^0\) is defined in \(U\) by the condition \(b\ne0\).

We have thus constructed a smooth affine \(S\)-group \(G^0\), with
connected fibers, possessing a self-centralizing maximal torus \(T\)
whose Weyl group \(W_{G^0}(T)\) is not finite \textup{(compare
\SGARefLink{sga3:xix:theorem:2:5}{Theorem 2.5})}.
\SGARefTarget{sga3:xix:paragraph:5:13}

\SGANextSectionTarget{sga3:xix:section:6}\SGASection{6}{Local existence of maximal tori.  The Weyl group}

In the proof of \SGARefLink{sga3:xix:theorem:2:5}{Theorem 2.5} we used a result from Exposé XI on the
local existence, for the étale topology, of maximal tori.  The proof
in Exposé XI uses a deep representability result \textup{(\SGARefLink{sga3:xi:theorem:4:1}{XI 4.1})}.
In the particular case under consideration, one can give another,
far more elementary proof, based on the ideas of Exposé XII, no.~7.

\par\medskip\noindent
\SGAInlineRefTarget{sga3:xix:proposition:6-1}\textbf{\textit{\SGARefLink{sga3:xix:proposition:6:1}{Proposition 6.1}.}}---
\textit{Let \(S\) be a scheme, let \(G\) be a smooth affine
\(S\)-group scheme with connected fibers over \(S\), and let
\(s_0\in S\) be a point such that the maximal tori of the geometric
fiber \(G_{\bar s_0}\) are self-centralizing.  There exist an étale
morphism \(S'\to S\) whose image contains \(s_0\) and a split maximal
torus \(T\) of \(G_{S'}\).}
\SGARefTarget{sga3:xix:proposition:6:1}

\SGAProof{Proof}
First, we may suppose that \(S\) is affine.
\footnote{Editor's note: the preceding sentence and the reference
to EGA IV below have been added.}
Since \(G\) is of finite presentation over \(S\), we may successively
suppose that \(S\) is noetherian, local, and then henselian with
separably closed residue field \textup{(cf.\ EGA IV, 8.12,
\S\S 8.8 and 18.8)}.  Thus let
\[
  S=\operatorname{Spec}(A),
  \qquad
  k=\kappa(s_0),
\]
where \(A\) is henselian and \(k\) is separably closed.  Choose a
maximal torus \(T_0\) of \(G_0=G_k\).  Such a torus exists, for
example because the scheme of maximal tori of \(G_0\) is smooth over
\(k\) \textup{(\SGARefLink{sga3:xii:7.1:item:c}{Exposé XII, 7.1(c)})}.  Since \(k\) is separably
closed, \(T_0\) is split \textup{(cf.\ \SGARefLink{sga3:x:proposition:1:4}{Exposé X, 1.4})}, and is
therefore given by a group monomorphism
\[
  f_0\colon \mathbb G_{m,k}^{\,r}\longrightarrow G_0.
\]

% END inputs/XIX/tex/components/23_expose_XIX_paragraph5_11_completion_through_proposition6_1_opening_page023.tex

% BEGIN inputs/XIX/tex/components/24_expose_XIX_proposition6_1_proof_through_proposition6_3_opening_page024.tex
Let \(m>1\) be an integer prime to the characteristic of \(k\).
By \SGARefLink{sga3:viii:corollary:6:7}{Exposé VIII, 6.7}, for every \(h>0\),
\(\operatorname{Centr}_{G_0}({}_{m^h}T_0)\) is represented by a
closed subscheme of \(G_0\).  The subgroup schemes
\({}_{m^h}T_0\) are schematically dense in \(T_0\)
\textup{(cf.\ \SGARefLink{sga3:ix:remark:4:10}{Exposé IX, 4.10})}, and \(G_0\) is noetherian.
Consequently there is an \(h\) such that
\[
  \operatorname{Centr}_{G_0}({}_{m^h}T_0)
  =
  \operatorname{Centr}_{G_0}(T_0)
  =
  T_0.
\]
Put \(n=m^h\).  Since \(n\) is invertible on \(S\),
\({}_n\mathbb G_{m,S}\) is isomorphic to
\((\mathbb Z/n\mathbb Z)_S\), so \(f_0\) defines a group
monomorphism
\[
  u_0\colon(\mathbb Z/n\mathbb Z)_k^{\,r}\longrightarrow G_0
\]
such that \(\operatorname{Centr}_{G_0}(u_0)=T_0\).

Now the \(S\)-functor
\[
  P=
  \operatorname{Hom}_{S\text{-}\mathrm{gr}}
  \bigl((\mathbb Z/n\mathbb Z)_S^{\,r},G\bigr)
\]
is represented by an \(S\)-scheme of finite type: it is a closed
subscheme of \(r\) copies of
\[
  {}_nG
  =
  \operatorname{Hom}_{S\text{-}\mathrm{gr}}
  \bigl((\mathbb Z/n\mathbb Z)_S,G\bigr)
  =
  \operatorname{Ker}\bigl(G\xrightarrow{\,n\,}G\bigr).
\]
But \(P\) is smooth over \(S\) \textup{(\SGARefLink{sga3:ix:theorem:3:6}{Exposé IX, 3.6})}.
Therefore \(u_0\in P(k)\) lifts to a section \(u\in P(S)\), by
Hensel's lemma \textup{(\SGARefLink{sga3:xi:corollary:1:11}{Exposé XI, 1.11})}:
\[
  u\colon(\mathbb Z/n\mathbb Z)_S^{\,r}\longrightarrow G.
\]

Consider \(H=\operatorname{Centr}_G(u)\).  It is a closed subgroup
scheme of \(G\), by \SGARefLink{sga3:viii:examples:6:5:item:e:01}{Exposé VIII, 6.5(e)}, and its special fiber is
\(H_0=T_0\).  Moreover, \(H\) is smooth over \(S\).  Indeed, let
\(S'=\operatorname{Spec}(A')\) be an affine \(S\)-scheme, let
\[
  u'\colon({}_n\mathbb G_{m,S'})^r\longrightarrow G_{S'}
\]
be obtained from \(u\) by base change, put
\(S'_J=\operatorname{Spec}(A'/J)\), where \(J\) is a square-zero
ideal, and let \(x\in H(S'_J)\).  Since \(G\) is smooth, \(x\) lifts
to an element \(g\in G(S')\).  Then
\[
  v=\operatorname{int}(g)(u')
\]
satisfies \(v_J=u'_J\).  By \SGARefLink{sga3:ix:theorem:3:2}{Exposé IX, 3.2}, there is therefore an
element \(g'\in G(S')\) such that \(g'_J=e\) and
\(\operatorname{int}(g')(v)=u'\).  It follows that \(h=g'g\) belongs
to \(H(S')\) and satisfies \(h_J=x\).
\footnote{Editor's note: the reference to Exposé VIII, 6.5(e) was
added above, and the preceding smoothness argument has been
expanded.}

Let \(H^0\) be the identity component of \(H\).  It is a subgroup
scheme of \(G\), smooth and with connected fibers, whose special
fiber is a torus.  By \SGARefLink{sga3:x:theorem:8-1}{Exposé X, 8.1}, it is a torus, necessarily split
by \SGARefLink{sga3:x:corollary:4-6}{Exposé X, 4.6}.  Put \(T=H^0\), and let
\[
  C=\operatorname{Centr}_G(T).
\]
This is a closed subgroup of \(G\) \textup{(\SGARefLink{sga3:viii:examples:6:5:item:e:01}{Exposé VIII, 6.5(e)})}
and is smooth \textup{(\SGARefLink{sga3:xi:corollary:2:4}{Exposé XI, 2.4})}.  Consider \(C^0\)
\textup{(in fact \(C^0=C\), but we do not need this)}.  Then
\(C^0\supset T\), and both are smooth groups with connected fibers.
They coincide at \(s_0\), hence in a neighborhood of \(s_0\).
After shrinking \(S\), we may therefore suppose that \(C^0=T\);
then, a fortiori, \(T\) is maximal.

\par\medskip\noindent
\SGAInlineRefTarget{sga3:xix:remark:6-2}\textbf{\textit{\SGARefLink{sga3:xix:remark:6:2}{Remark 6.2}.}}---
\textit{The proof shows in particular that the reductive rank of
\(G_{\bar s}\) is constant in a neighborhood of \(s=s_0\).}
\SGARefTarget{sga3:xix:remark:6:2}

\par\medskip\noindent
\SGAInlineRefTarget{sga3:xix:proposition:6-3}\textbf{\textit{\SGARefLink{sga3:xix:proposition:6:3}{Proposition 6.3}.}}---
\textit{Let \(S\) be a scheme, let \(G\) be an \(S\)-group scheme
that is smooth and of finite presentation over \(S\), and let \(Q\)
be a subtorus of \(G\).}
\SGARefTarget{sga3:xix:proposition:6:3}

% END inputs/XIX/tex/components/24_expose_XIX_proposition6_1_proof_through_proposition6_3_opening_page024.tex

% BEGIN inputs/XIX/tex/components/25_expose_XIX_proposition6_3_completion_and_bibliography_page025.tex
\begin{enumerate}
\item[\textup{(i)}]
\textit{\(\operatorname{Centr}_G(Q)\) and
\(\operatorname{Norm}_G(Q)\) are represented by closed subgroup
schemes, smooth \textup{(and hence of finite presentation)} over
\(S\).}

\item[\textup{(ii)}]
\textit{\(\operatorname{Centr}_G(Q)\) is an open and closed
subscheme of \(\operatorname{Norm}_G(Q)\).  The quotient}
\[
  W_G(Q)
  =
  \operatorname{Norm}_G(Q)/\operatorname{Centr}_G(Q)
\]
\textit{is represented by an open subgroup scheme of
\(\operatorname{Aut}_{S\text{-}\mathrm{gr}}(Q)\); it is therefore a
quasi-finite, étale, separated \(S\)-group scheme.}

\item[\textup{(iii)}]
\textit{For every \(s\in S\), put}
\[
  w(s)
  =
  \left|
  \operatorname{Norm}_{G(\bar s)}(Q(\bar s))/
  \operatorname{Centr}_{G(\bar s)}(Q(\bar s))
  \right|.
\]
\textit{Then \(s\longmapsto w(s)\) is lower semicontinuous, and is
constant in a neighborhood of \(s\) if and only if \(W_G(Q)\) is
finite over \(S\) in a neighborhood of \(s\).}
\end{enumerate}

\SGAProof{Proof}
By \SGARefLink{sga3:xi:corollary:6:11}{Exposé XI, 6.11},
\(\operatorname{Centr}_G(Q)\) and \(\operatorname{Norm}_G(Q)\) are
represented by closed subschemes of \(G\) of finite presentation.
They are smooth by \SGARefLink{sga3:xi:corollary:2:4}{Exposé XI, 2.4} and 2.4 bis, which proves
\textup{(i)}.  Assertions \textup{(ii)} and \textup{(iii)} are then
proved as in \SGARefLink{sga3:xi:proposition:5:9}{Exposé XI, 5.9} and 5.10; those proofs in fact use only
\textup{(i)}, and not the deeper theorems of \SGARefLink{sga3:xi:theorem:4:1}{Exposé XI, 4.1} and 4.2.

\section*{Bibliography}
\addcontentsline{toc}{section}{Bibliography}

\begin{sloppypar}
\begin{description}
\item[{[Bible]}]
C.~Chevalley, with the collaboration of P.~Cartier,
A.~Grothendieck, and M.~Lazard,
\textit{Classification des groupes de Lie algébriques}, 1956--58.

\item[{[Ch05]}]
C.~Chevalley,
\textit{Classification des groupes algébriques semi-simples}, with
the collaboration of P.~Cartier, A.~Grothendieck, and M.~Lazard,
\textit{Collected Works}, vol.~3, Springer, 2005.

\item[{[Tô55]}]
C.~Chevalley, \textit{Sur certains groupes simples},
\textit{Tôhoku Math.\ J.}\ (2) \textbf{7} (1955), 14--66.

\item[{[TO70]}]
J.~Tate and F.~Oort, \textit{Group schemes of prime order},
\textit{Ann.\ scient.\ Éc.\ Norm.\ Sup.}\ (4), t.~\textbf{3}
(1970), 1--21.
\end{description}
\end{sloppypar}

% END inputs/XIX/tex/components/25_expose_XIX_proposition6_3_completion_and_bibliography_page025.tex


